% Computer Science A --- Angular (framework). Angular Introduction unit.

\runningheader{Computer Science A}{Web: Angular}{Page \thepage\ of \numpages}

\begingradingrange{csawebangular}

\uplevel{\par\textbf{Part 1: Angular introduction}\par}

\question[4]
Angular is a front-end framework developed primarily by:
\begin{choices}
  \choice Microsoft
  \CorrectChoice Google
  \choice Twitter
  \choice Facebook
\end{choices}
\begin{solution}
  Google maintains modern Angular. Other companies (Microsoft, Meta, etc.) build products \emph{with} it.
\end{solution}

\question[4]
Which statement best describes modern Angular?
\begin{choices}
  \choice A markup language that replaces HTML
  \choice A relational database for web servers
  \CorrectChoice A TypeScript-centered, component-based \textbf{framework} for scalable SPAs (and often PWAs)
  \choice A CSS preprocessor only
\end{choices}
\begin{solution}
  Angular organizes apps into \textbf{components}, ships libraries (routing, HTTP), and includes CLI tooling. \textbf{SPAs} load once and swap views; \textbf{PWAs} add app-like delivery (installable, responsive). It structures HTML/CSS/TS---it does not replace them.
\end{solution}

\question[4]
Which statement about Angular's language history is most accurate?
\begin{choices}
  \choice Angular has always been TypeScript-only with no JavaScript heritage
  \CorrectChoice AngularJS (\texttt{1.x}) used JavaScript (\texttt{2010}); modern Angular (\texttt{2+}, \texttt{2016}) centers on TypeScript
  \choice \texttt{Angular~2} replaced TypeScript with Python
  \choice Angular templates compile to SQL
\end{choices}
\begin{solution}
  Do not confuse \textbf{AngularJS} (\texttt{1.x}, JS, MVC, directives) with \textbf{Angular} (\texttt{2+}, TypeScript, components, CLI).
\end{solution}

\question[4]
Before starting Angular in this course, which background is expected?
\begin{choices}
  \choice Only SQL and shell scripting
  \CorrectChoice HTML, CSS, and JavaScript, plus basic MVC and TypeScript ideas
  \choice Only AngularJS \texttt{1.x} with no web fundamentals
  \choice Server-side Java only
\end{choices}
\begin{solution}
  Prerequisites: sound HTML/CSS/JS, basic \textbf{MVC}, and introductory \textbf{TypeScript}.
\end{solution}

\question[4]
Which description best matches a \textbf{Single Page Application (SPA)}?
\begin{choices}
  \choice A site that reloads a full HTML document on every click by default
  \CorrectChoice A web app loaded once that updates views with components instead of constant full page reloads
  \choice A mobile app installed only from an app store
  \choice A spreadsheet in a SQL table
\end{choices}
\begin{solution}
  SPAs load once, then update via client-side routing and components---faster perceived navigation than many MPAs.
\end{solution}

\question[4]
According to the introduction unit, Angular is especially appropriate when:
\begin{choices}
  \choice You need a one-screen static handout with no interactivity
  \CorrectChoice You are building large-scale, dynamic applications with reusable components and many developers
  \choice You must avoid any client--server communication
  \choice You want to replace CSS with SQL
\end{choices}
\begin{solution}
  Prefer Angular for large dynamic apps, SPAs/PWAs, cross-platform browser targets, and teams. Compared with a small utility \textbf{library}, a \textbf{framework} adds opinionated layout, official libraries, and integrated build/test/deploy tooling.
\end{solution}

\question[4]
Which is \textbf{not} a valid reason to choose Angular listed in the introduction unit?
\begin{choices}
  \choice Building cross-platform web experiences
  \choice Organizing a large project with reusable components
  \CorrectChoice Eliminating the need for HTML, CSS, and JavaScript entirely
  \choice Creating SPA or PWA experiences
\end{choices}
\begin{solution}
  Angular scales the web stack; it does not remove HTML templates, CSS, or TypeScript/JavaScript logic.
\end{solution}

\newpage

\uplevel{\par\textbf{Part 2: \texttt{Angular~2} vs AngularJS and project setup}\par}

\question[4]
When moving from \textbf{AngularJS} (\texttt{1.x}) to modern \textbf{Angular} (\texttt{2+}), which mindset is most accurate?
\begin{choices}
  \choice Every AngularJS pattern transfers unchanged
  \CorrectChoice Treat \texttt{2+} as a fresh mental model---components, modules, and TypeScript replace controllers, \texttt{\$scope}, and many \texttt{1.x} idioms
  \choice \texttt{2+} removes components entirely
  \choice \texttt{2+} cannot use TypeScript
\end{choices}
\begin{solution}
  \texttt{2+} uses \textbf{components} (class + template) in \textbf{modules}, with decorators and TypeScript imports/exports:
\begin{lstlisting}[style=tsexam]
@Component({
  selector: 'app-root',
  template: `<h1>{{ title }}</h1>`,
})
export class AppComponent {
  title = 'Hello';
}
\end{lstlisting}
  Comfortable TypeScript features include modules, interfaces, classes, and decorators.
\end{solution}

\question[4]
In \texttt{2+}, the \textbf{model} and \textbf{view} are most commonly defined where?
\begin{choices}
  \choice Model in CSS; view in \texttt{package.json}
  \CorrectChoice Model in a TypeScript class; view in the component template (inline or separate HTML)
  \choice Both only in SQL
  \choice Both only in \texttt{localStorage}
\end{choices}
\begin{solution}
  The \textbf{model} is usually a TypeScript class; the \textbf{view} is the component template (HTML). Logic stays in \lstinline[style=tsinline]{.ts} files; presentation stays in templates.
\end{solution}

\question[4]
To install tooling for local Angular development, which setup is correct?
\begin{choices}
  \choice \texttt{javac} and \texttt{jar} only
  \CorrectChoice \texttt{Node.js} and \texttt{npm} (Node runs JS; npm manages packages), then project scripts or the CLI
  \choice \texttt{python} and \texttt{pip} only
  \choice No runtime; edit files only in the browser
\end{choices}
\begin{solution}
  Verify the toolchain:
\begin{lstlisting}[style=cliexam]
node -v
npm -v
\end{lstlisting}
  Node runs JavaScript tooling; npm installs packages. Use current LTS versions in practice.
\end{solution}

\question[4]
Inside the project folder, \texttt{npm start} in the quickstart workflow typically:
\begin{choices}
  \choice Deletes \texttt{src/} permanently
  \CorrectChoice Starts a dev server (often \texttt{localhost:3000}), watches files, and recompiles TypeScript before the UI updates
  \choice Pushes to GitHub automatically
  \choice Installs the CLI globally
\end{choices}
\begin{solution}
  Saving edits (e.g.\ in the root component) triggers transpile output; the browser shows updated bindings after the build step.
\end{solution}

\question[4]
Which trio are the main startup files under \texttt{src/app/} in the quickstart layout?
\begin{choices}
  \choice \texttt{package.json}, \texttt{README.md}, \texttt{.gitignore}
  \CorrectChoice \texttt{main.ts}, \texttt{app-module.ts}, and \texttt{app-component.ts}
  \choice \texttt{main.ts}, \texttt{\$scope.js}, and \texttt{controller.ts}
  \choice \texttt{index.css}, \texttt{jquery.js}, \texttt{bootstrap.min.js}
\end{choices}
\begin{solution}
  Startup layout under \lstinline[style=cliinline]{src/app/}:
\begin{lstlisting}[style=tsexam]
// main.ts        -- bootstraps AppModule
// app.module.ts  -- declares/imports components
// app.component.ts -- root UI component
\end{lstlisting}
  New features go in \lstinline[style=cliinline]{src/app/}; project-root files configure build and tests. \lstinline[style=tsinline]{AppModule} organizes the app---it does not transpile TypeScript (the toolchain does).
\end{solution}

\question[4]
Which is \textbf{not} a responsibility of the root \texttt{AppModule}?
\begin{choices}
  \choice Declaring or importing application components
  \choice Linking the root component into the module
  \CorrectChoice Transpiling TypeScript inside the browser tab with no build step
  \choice Acting as the top-level module for the app
\end{choices}
\begin{solution}
  \lstinline[style=tsinline]{AppModule} declares which components belong to the app. Compilation is handled by \lstinline[style=cliinline]{npm start}, \lstinline[style=cliinline]{ng serve}, or the build bundler---not by the module class at runtime.
\end{solution}

\newpage

\uplevel{\par\textbf{Part 3: Angular environment setup}\par}

\question[4]
According to the Angular setup unit, a complete local toolchain requires (in order):
\begin{choices}
  \choice Angular CLI only, with no Node or npm
  \choice SQL Server, then Git
  \CorrectChoice Node.js from \texttt{nodejs.org}, npm (bundled with Node), then \texttt{npm install -g @angular/cli}
  \choice A browser plugin that eliminates TypeScript
\end{choices}
\begin{solution}
  Install in order: Node from \texttt{nodejs.org}, then the CLI globally:
\begin{lstlisting}[style=cliexam]
npm install -g @angular/cli
\end{lstlisting}
  \textbf{npm} (Node Package Manager) ships with Node. The \lstinline[style=cliinline]{-g} flag puts \lstinline[style=cliinline]{ng} on your PATH. On Windows, PowerShell execution policies may block npm scripts.
\end{solution}

\question[4]
Evaluate the statements (Angular setup unit):

\textbf{Statement 1:} Modern Angular projects are developed using npm and the Angular CLI.

\textbf{Statement 2:} AngularJS (\texttt{1.x}) projects use \emph{only} a third-party IDE, with no role for npm or package managers.
\begin{choices}
  \choice Both statements are true
  \CorrectChoice Statement 1 is true; Statement 2 is false
  \choice Statement 1 is false; Statement 2 is true
  \choice Both statements are false
\end{choices}
\begin{solution}
  Modern workflow: Node $\rightarrow$ npm $\rightarrow$ global CLI $\rightarrow$ \lstinline[style=cliinline]{ng}. AngularJS-era projects could use editors and other build tools; ``only an IDE, never npm'' is false.
\end{solution}

\question[4]
Which item is \textbf{not} listed as a setup requirement?
\begin{choices}
  \choice Node.js
  \choice npm
  \CorrectChoice A built-in Oracle database installed before Node
  \choice Angular CLI installed through npm
\end{choices}
\begin{solution}
  Setup needs Node, npm, and the CLI---not a proprietary database bundled with Angular.
\end{solution}

\newpage

\uplevel{\par\textbf{Part 4: AngularJS (\texttt{1.x}) core architecture}\par}

\question[4]
In the AngularJS basics unit, \textbf{AngularJS} is described as:
\begin{choices}
  \choice A static HTML validator with no JavaScript
  \CorrectChoice A JavaScript framework for dynamic websites with cross-browser support, extensibility, and two-way data binding
  \choice A relational database server
  \choice Identical to modern \texttt{Angular~14}
\end{choices}
\begin{solution}
  AngularJS (\texttt{1.x}) runs in the browser. \textbf{Two-way binding} syncs view and model automatically. ``Extensible'' means custom directives, services, and modules. Read context when a lesson says ``Angular''---it may mean AngularJS.
\end{solution}

\question[4]
Which set lists the \textbf{core building blocks} in the AngularJS overview?
\begin{choices}
  \choice Decorators, RxJS, and the CLI only
  \CorrectChoice Directives, Controllers, Modules, Services, Expressions and Filters, and Routing
  \choice \texttt{main.ts}, \texttt{app-module.ts}, \texttt{app-component.ts}
  \choice JVM, JDBC, and Servlets
\end{choices}
\begin{solution}
  \textbf{Directives} extend HTML; \textbf{controllers} connect view and model (often via \lstinline[style=tsinline]{\$scope}); \textbf{modules} partition the app; \textbf{services} share logic app-wide. Template expressions and filters format data, e.g.:
\begin{lstlisting}[style=jsexam]
{{ name | uppercase }}
\end{lstlisting}
  \textbf{Routing} maps URLs to templates for SPA-style navigation without full reloads.
\end{solution}

\question[4]
Which pairing correctly matches an AngularJS building block to its responsibility?
\begin{choices}
  \choice Directives $\rightarrow$ share HTTP logic app-wide; Services $\rightarrow$ bind HTML attributes
  \CorrectChoice Directives $\rightarrow$ extend HTML with behavior; Services $\rightarrow$ reusable app-wide logic
  \choice Controllers $\rightarrow$ npm registry; Modules $\rightarrow$ V8 engine
  \choice Routing $\rightarrow$ uppercase filter only; Filters $\rightarrow$ URL mapping only
\end{choices}
\begin{solution}
  Services are \emph{shared} across the app---not limited to one HTML tag.
\end{solution}

\question[4]
When studying \textbf{modern} Angular (\texttt{2+}) after AngularJS lessons, you should:
\begin{choices}
  \choice Assume every ``Angular'' sentence means AngularJS \texttt{1.x}
  \CorrectChoice Read context: AngularJS uses controllers/directives/MVC; modern Angular uses TypeScript components, modules, and the CLI
  \choice Use \texttt{\$scope} in every new CLI project
  \choice Ignore TypeScript because both versions are identical
\end{choices}
\begin{solution}
  Naming overlap is historical. Later parts describe modern Angular unless the question says AngularJS.
\end{solution}

\newpage

\uplevel{\par\textbf{Part 5: Angular CLI}\par}

\question[4]
The \textbf{Angular CLI} is:
\begin{choices}
  \choice A browser extension that edits CSS without a build step
  \CorrectChoice A command-line tool (\texttt{ng}) used to scaffold, build, and serve Angular applications
  \choice A replacement for HTML in every template
  \choice The same thing as AngularJS \texttt{\$scope}
\end{choices}
\begin{solution}
  \lstinline[style=cliinline]{ng} stands for \textbf{Angular}. Install with \lstinline[style=cliinline]{npm install -g @angular/cli}. The CLI standardizes scaffold/build/serve workflows. CLI major versions align with Angular major versions; minor CLI releases may ship separately.
\end{solution}

\question[4]
Which command creates a \textbf{new} Angular project named \texttt{my-app}?
\begin{choices}
  \choice \texttt{ng create my-app}
  \CorrectChoice \texttt{ng new my-app}
  \choice \texttt{ng build my-app}
  \choice \texttt{ng serve my-app}
\end{choices}
\begin{solution}
  \lstinline[style=cliinline]{cd} to the parent folder, then scaffold:
\begin{lstlisting}[style=cliexam]
ng new my-app
\end{lstlisting}
  \lstinline[style=cliinline]{ng serve} runs a dev server for an \emph{existing} project---not for creating one.
\end{solution}

\question[4]
Which command starts a local dev server for a CLI project on port \textbf{\texttt{4200}} by default?
\begin{choices}
  \choice \texttt{ng new}
  \CorrectChoice \texttt{ng serve}
  \choice \texttt{ng help} only
  \choice \texttt{npm install -g @angular/cli} only
\end{choices}
\begin{solution}
\begin{lstlisting}[style=cliexam]
ng serve
ng serve --open
\end{lstlisting}
  Default URL: \texttt{http://localhost:4200/}. \lstinline[style=cliinline]{--open} launches the browser. \lstinline[style=cliinline]{ng help} lists commands. (Older quickstarts used \lstinline[style=cliinline]{npm start} on port \texttt{3000}.)
\end{solution}

\question[4]
The learning objective for the Angular CLI unit is to:
\begin{choices}
  \choice Memorize every Forbes article from 1990
  \CorrectChoice Create an Angular project using the Angular CLI
  \choice Install AngularJS controllers only
  \choice Avoid npm entirely
\end{choices}
\begin{solution}
  Outcome: install CLI, then \lstinline[style=cliinline]{ng new} and \lstinline[style=cliinline]{ng serve}. Modern Angular projects are commonly CLI-driven.
\end{solution}

\newpage

\uplevel{\par\textbf{Part 6: Single Page Applications}\par}

\question[4]
A \textbf{Single Page Application (SPA)} is best defined as a web application that:
\begin{choices}
  \choice Reloads a full HTML document from the server on every click by default
  \CorrectChoice Dynamically rewrites a single web page with new data from the server
  \choice Stores all content only in a database with no browser UI
  \choice Runs only as a native mobile binary
\end{choices}
\begin{solution}
  First request returns HTML (the shell); later updates use JavaScript APIs and server data without a full document reload each time.
\end{solution}

\question[4]
Compared with a traditional \textbf{multi-page application (MPA)}, SPAs often provide:
\begin{choices}
  \choice Slower navigation because HTML reloads every time
  \CorrectChoice Quicker subsequent navigation, app-like UX, easier feature growth, and less repeated HTML download (often JSON updates instead)
  \choice No use of JavaScript
  \choice Guaranteed better SEO than every static site
\end{choices}
\begin{solution}
  Benefits: load shell once, fluid UX, extensible components/routes, lower repeated bandwidth. Drawbacks: SEO challenges for pure client rendering, heavy browser memory/CPU, and XSS risk if DOM updates are unsafe.
\end{solution}

\question[4]
Which statement about SPA is \textbf{false}?
\begin{choices}
  \choice SPAs can offer quick loading after the initial shell is cached
  \choice SPAs can use less bandwidth than repeatedly downloading full HTML pages
  \CorrectChoice SPAs typically use \emph{fewer} browser resources than minimal static pages
  \choice SPAs can provide a seamless, app-like user experience
\end{choices}
\begin{solution}
  The false claim is ``fewer browser resources.'' SPAs keep large scripts, routing state, and DOM in memory after the first load, so they usually demand \emph{more} CPU and memory than a minimal static page---not less. Quick loading, lower repeated bandwidth, and app-like UX are true benefits.
\end{solution} 

\question[4]
Angular fits SPA development because it provides:
\begin{choices}
  \choice Only static \texttt{.txt} files
  \CorrectChoice Components, client-side routing, and CLI tooling for single-load, dynamically updated apps
  \choice Mandatory full page reload on every route
  \choice A ban on TypeScript
\end{choices}
\begin{solution}
  Core SPA pattern: one load, then dynamic updates---matching Angular's component and router model.
\end{solution}

\newpage

\uplevel{\par\textbf{Part 7: Webpack and Angular build systems}\par}

\question[4]
\textbf{Webpack} is best described as:
\begin{choices}
  \choice A relational database for Angular components
  \CorrectChoice A module bundler that builds optimized browser bundles from application source
  \choice A CSS-only preprocessor
  \choice The same executable as \texttt{ng}
\end{choices}
\begin{solution}
  A \textbf{bundle} groups related files (JS, CSS, HTML, etc.) for one client response. Config concepts: \textbf{entry} (where the graph starts), \textbf{output} (where bundles are written), \textbf{loaders} (non-JS file types), \textbf{plugins} (objects with \texttt{apply} for lifecycle hooks). Angular, React, and Vue historically used Webpack for bundling.
\end{solution}

\question[4]
When you run \texttt{ng new} or \texttt{ng serve}, what changed between older (Webpack-based) and modern Angular?
\begin{choices}
  \choice The commands were renamed to \texttt{webpack new} and \texttt{webpack serve}
  \CorrectChoice The commands stayed the same; the internal build engine changed---modern Angular uses a faster framework-managed system without exposing \texttt{webpack.config.js}
  \choice Modern Angular cannot use TypeScript
  \choice Older Angular never used npm
\end{choices}
\begin{solution}
  Older CLI (through \texttt{Angular~17}) managed Webpack internally. \texttt{ng serve} used Webpack Dev Server; modern \texttt{ng serve} uses Angular's optimized dev server.
\end{solution}

\question[4]
\texttt{ng eject} historically allowed developers to expose \texttt{webpack.config.js}. Why was it removed?
\begin{choices}
  \choice Ejected projects upgraded automatically with every CLI release
  \CorrectChoice Ejected projects were hard to maintain, blocked CLI improvements, and left fragile configs
  \choice Webpack stopped supporting TypeScript
  \choice Modern Angular forbids all bundling
\end{choices}
\begin{solution}
  Ejecting gave full Webpack control (entry, loaders, \texttt{devServer} on port \texttt{4200}, \texttt{historyApiFallback} for SPA routes) but transferred full maintenance burden to the team.
\end{solution}

\question[4]
Modern Angular projects configure builds primarily through:
\begin{choices}
  \choice A required root \texttt{webpack.config.js} in every repo
  \CorrectChoice \texttt{angular.json}---for example \texttt{main}, \texttt{index}, \texttt{outputPath}, \texttt{assets}, \texttt{styles}
  \choice Only inline \texttt{<style>} tags
  \choice The Windows Registry
\end{choices}
\begin{solution}
  Modern builds use \lstinline[style=cliinline]{angular.json} instead of hand-editing \lstinline[style=cliinline]{webpack.config.js}:
\begin{lstlisting}[style=jsonexam]
"build": {
  "options": {
    "main": "src/main.ts",
    "outputPath": "dist/my-app",
    "assets": ["src/assets"],
    "styles": ["src/styles.css"]
  }
}
\end{lstlisting}
  \lstinline[style=cliinline]{ng build} minifies, splits bundles, and hashes filenames by default.
\end{solution}

\question[4]
Which statement about Webpack and modern Angular is most accurate?
\begin{choices}
  \choice Every new project must eject before \texttt{ng build}
  \CorrectChoice Modern Angular still bundles and optimizes, but developers usually configure \texttt{angular.json}, not raw Webpack
  \choice The CLI never used Webpack
  \choice \texttt{webpack.config.js} is required in all \texttt{Angular~18+} apps
\end{choices}
\begin{solution}
  Webpack knowledge matters for \texttt{Angular~17} and earlier legacy codebases; newer projects abstract the bundler while still producing optimized output.
\end{solution}

\newpage

\uplevel{\par\textbf{Part 8: Angular version history}\par}

\question[4]
\textbf{AngularJS} (\texttt{1.x}) versus \textbf{Angular}~\texttt{2+} differ mainly in that:
\begin{choices}
  \choice Both are TypeScript-only frameworks released in \texttt{2016}
  \CorrectChoice AngularJS is JavaScript-based (\texttt{2010}, \texttt{1.7.x}, LTS), MVC/directives; \texttt{Angular~2+} is TypeScript-based (\texttt{2016}), component-based, mobile-oriented, CLI-driven
  \choice Only AngularJS provides a CLI
  \choice Modern Angular uses \texttt{\$scope} instead of components
\end{choices}
\begin{solution}
  AngularJS (\texttt{1.x}): JavaScript, MVC, directives like \lstinline[style=jsexam]{ng-model}. Modern Angular (\texttt{2+}): TypeScript, components, CLI:
\begin{lstlisting}[style=tsexam]
// modern template binding
<input [(ngModel)]="name">
\end{lstlisting}
  Only \texttt{1.x} used JavaScript as the framework language.
\end{solution}

\question[4]
Why was \texttt{Angular~3} skipped?
\begin{choices}
  \choice TypeScript was removed
  \CorrectChoice \texttt{@angular/router} would have been \texttt{4.x} while \texttt{@angular/core} and \texttt{@angular/compiler} were still \texttt{2.x}---the public release jumped to \texttt{Angular~4} to avoid confusion
  \choice \texttt{Angular~3} merged into AngularJS \texttt{1.8}
  \choice The CLI forbade version numbers above \texttt{2}
\end{choices}
\begin{solution}
  \texttt{Angular~4} (\texttt{2017}) had relatively few changes versus \texttt{Angular~2}---mostly package version alignment.
\end{solution}

\question[4]
Which version introduced the \textbf{Ivy} compiler and runtime by default?
\begin{choices}
  \choice \texttt{Angular~5}
  \choice \texttt{Angular~7}
  \CorrectChoice \texttt{Angular~9}
  \choice \texttt{Angular~14}
\end{choices}
\begin{solution}
  \textbf{Answer: \texttt{Angular~9}.}

  \textbf{Ivy} is Angular's default \textbf{compilation and rendering engine} (from \texttt{Angular~9} onward). It replaced the older \textbf{View Engine}. Ivy compiles templates and components into instructions the browser can run more directly, which helps reduce bundle size, improve build performance, and make debugging clearer than the previous pipeline.

  Other notable major releases (organized for review):
  \begin{itemize}
    \item \texttt{Angular~5}: Build Optimizer, Universal state transfer, compiler improvements, i18n pipes
    \item \texttt{Angular~6}: Angular Elements, CDK, CLI workspaces, schematics
    \item \texttt{Angular~7}: CLI prompts, Material/CDK updates, virtual scrolling, drag and drop
    \item \texttt{Angular~8}: differential loading, CLI builder APIs, web worker support
    \item \texttt{Angular~12}: strict mode by default; HTTP interceptor context
    \item \texttt{Angular~14}: standalone components, strictly typed forms, CLI auto-completion
  \end{itemize}
\end{solution}

\question[4]
Which version highlights include \textbf{standalone components} and strictly typed forms?
\begin{choices}
  \choice \texttt{Angular~9}
  \choice \texttt{Angular~12}
  \CorrectChoice \texttt{Angular~14}
  \choice AngularJS \texttt{1.7}
\end{choices}
\begin{solution}
  \texttt{v14} also added CLI auto-completion. \texttt{v13} simplified dynamic components; \texttt{v11} added automatic font inlining---see \texttt{angular.io/guide/releases}.
\end{solution}

\question[4]
The version history summary states that Angular:
\begin{choices}
  \choice Started in \texttt{2016} with TypeScript only
  \CorrectChoice Started in \texttt{2010} as JavaScript (AngularJS) and shifted to TypeScript from \texttt{version~2} (\texttt{2016}), with ongoing major releases
  \choice Stopped updating after \texttt{version~2}
  \choice Renamed TypeScript to AngularJS in \texttt{2017}
\end{choices}
\begin{solution}
  Timeline: \texttt{2010} JS (\texttt{1.x}) $\rightarrow$ \texttt{2016} TS rewrite (\texttt{2+}) $\rightarrow$ continued semver majors.
\end{solution}

\newpage

\uplevel{\par\textbf{Part 9: The \texttt{@Component} decorator}\par}

\question[4]
An Angular \textbf{component} is best described as:
\begin{choices}
  \choice A npm package that replaces TypeScript
  \CorrectChoice A building block of the application made of a template, selector, optional styles, and a TypeScript class
  \choice A SQL table used only on the server
  \choice The same thing as a Webpack loader
\end{choices}
\begin{solution}
  Components are the building blocks of Angular apps. Each component typically includes:
  \begin{itemize}
    \item an \textbf{HTML template} (view),
    \item a \textbf{CSS selector} (where the component is inserted),
    \item optional \textbf{CSS styles}, and
    \item a \textbf{TypeScript class} (behavior and data).
  \end{itemize}
  Components are independent, support lazy loading, and help maintain large SPAs.
\end{solution}

\question[4]
Which package defines the \texttt{@Component} decorator?
\begin{choices}
  \CorrectChoice \texttt{@angular/core}
  \choice \texttt{@angular/http}
  \choice \texttt{@angular/component}
  \choice \texttt{@angular/directives}
\end{choices}
\begin{solution}
  Import \lstinline[style=tsinline]{Component} from \lstinline[style=tsinline]{@angular/core}. The decorator marks a class as a component and supplies runtime metadata.
\end{solution}

\question[4]
In \texttt{@Component} metadata, \texttt{selector} specifies:
\begin{choices}
  \choice The path to the HTML template file
  \CorrectChoice The CSS selector where Angular inserts this component in a parent template
  \choice The npm registry name for the app
  \choice The port for \texttt{ng serve}
\end{choices}
\begin{solution}
  Example: \lstinline[style=tsinline]{selector: 'app-root'} means Angular creates this component wherever it finds \lstinline[style=tsinline]{<app-root></app-root>} in HTML.
\end{solution}

\question[4]
Which metadata property points to the HTML file for a component's view?
\begin{choices}
  \choice \texttt{animations}
  \choice \texttt{exportAs}
  \choice \texttt{providers}
  \CorrectChoice \texttt{templateUrl}
\end{choices}
\begin{solution}
  \texttt{templateUrl} is the module-relative path to an external HTML file (for example \lstinline[style=tsinline]{./app.component.html}). Alternatively, use inline \texttt{template} with the HTML in the decorator.
\end{solution}

\question[4]
The default root component in a new CLI project is typically split across:
\begin{choices}
  \choice \texttt{main.ts}, \texttt{package.json}, and \texttt{README.md} only
  \CorrectChoice \texttt{app.component.ts}, \texttt{app.component.html}, and \texttt{app.component.css}
  \choice \texttt{webpack.config.js} only
  \choice \texttt{angular.json} only
\end{choices}
\begin{solution}
\begin{lstlisting}[style=tsexam]
import { Component } from '@angular/core';

@Component({
  selector: 'app-root',
  templateUrl: './app.component.html',
  styleUrls: ['./app.component.css'],
})
export class AppComponent {
  title = 'my first app';
}
\end{lstlisting}
  \texttt{styleUrls} lists CSS files for the template. A component must belong to an \texttt{NgModule} to be used by other components.
\end{solution}

\question[4]
Which command creates a new component named \texttt{user-profile}?
\begin{choices}
  \choice \texttt{ng new user-profile}
  \CorrectChoice \texttt{ng generate component user-profile} or \texttt{ng g c user-profile}
  \choice \texttt{ng serve user-profile}
  \choice \texttt{npm install user-profile}
\end{choices}
\begin{solution}
  The CLI scaffolds a folder with \texttt{.component.ts}, \texttt{.component.html}, \texttt{.component.css}, and \texttt{.component.spec.ts}.
\end{solution}

\question[4]
\texttt{@Component} is a special form of the \texttt{@Directive} decorator.
\begin{choices}
  \CorrectChoice True
  \choice False
\end{choices}
\begin{solution}
  \textbf{True.} Components are directives with a template. You declare a component class by adding \texttt{@Component}---not by implementing a separate \texttt{Component} interface alone.
\end{solution}

\question[4]
Every Angular application has at least one component (the root component).
\begin{choices}
  \CorrectChoice True
  \choice False
\end{choices}
\begin{solution}
  \textbf{True.} Apps bootstrap through a root component (often \texttt{AppComponent}). You typically define one root \texttt{NgModule} and attach feature modules to it.
\end{solution}

\newpage

\uplevel{\par\textbf{Part 10: Templates}\par}

\question[4]
An Angular \textbf{template} is:
\begin{choices}
  \choice A JavaScript file that replaces TypeScript
  \CorrectChoice A blueprint for part of the UI, written in HTML with Angular syntax
  \choice The same as \texttt{package.json}
  \choice A Webpack configuration file
\end{choices}
\begin{solution}
  Templates describe what the user sees for a component. They are HTML-based (with Angular extensions) and tie directly to the component \textbf{view}. They are usually a \emph{fragment} of a page---not a full document with \lstinline[style=htmlexam]{<html>} or \lstinline[style=htmlexam]{<body>}.
\end{solution}

\question[4]
Which properties attach a template to a component?
\begin{choices}
  \choice \texttt{selector} only
  \choice \texttt{styleUrls} only
  \CorrectChoice \texttt{template} (inline) or \texttt{templateUrl} (external file)
  \choice \texttt{bootstrap} only
\end{choices}
\begin{solution}
  \textbf{Inline} template---HTML in backticks on \texttt{template}:
\begin{lstlisting}[style=tsexam]
@Component({
  selector: 'app-root',
  template: `<h1>App Component</h1>
             <p>Inline template</p>`,
})
export class AppComponent { }
\end{lstlisting}
  \textbf{External} template---separate file via \texttt{templateUrl}:
\begin{lstlisting}[style=tsexam]
@Component({
  selector: 'app-root',
  templateUrl: './app.component.html',
})
export class AppComponent { }
\end{lstlisting}
\end{solution}

\question[4]
Which elements are \textbf{not} supported inside Angular component templates?
\begin{choices}
  \choice \texttt{<div>} and \texttt{<p>}
  \CorrectChoice \texttt{<script>}, and document-level tags such as \texttt{<html>}, \texttt{<body>}, and \texttt{<base>}
  \choice \texttt{<button>} and \texttt{<input>}
  \choice \texttt{<h1>}--\texttt{<h6>} only
\end{choices}
\begin{solution}
  \texttt{<script>} is disallowed (keeps logic in the component class and reduces injection risk). Templates are injected into part of a page, so full-document tags like \lstinline[style=htmlexam]{<html>} and \lstinline[style=htmlexam]{<body>} are unnecessary.
\end{solution}

\question[4]
A best practice for external template and style files is to:
\begin{choices}
  \choice Use random unrelated file names in \texttt{node\_modules/}
  \CorrectChoice Name them to match the component (for example \texttt{app.component.html} for \texttt{AppComponent})
  \choice Store them only in the repository root with no path
  \choice Embed all CSS inside \texttt{main.ts}
\end{choices}
\begin{solution}
  Pair \texttt{app.component.ts} with \texttt{app.component.html} and \texttt{app.component.css}. Use a module-relative path in \texttt{templateUrl} (for example \lstinline[style=tsinline]{./app.component.html}).
\end{solution}

\newpage

\uplevel{\par\textbf{Part 11: Template statements}\par}

\question[4]
\textbf{Template statements} in Angular are:
\begin{choices}
  \choice SQL queries run on the server only
  \CorrectChoice Methods or logic in the template, invoked in response to user events (for example a button click)
  \choice The same as Webpack entry points
  \choice Only CSS class names
\end{choices}
\begin{solution}
  Template statements handle user actions (clicks, keystrokes, submits). They use a JavaScript-\emph{like} syntax, but the parser rules differ from template \textbf{expressions}.
\end{solution}

\question[4]
Which fragment is a valid template statement binding?
\begin{choices}
  \choice \texttt{[value]="count++"}
  \CorrectChoice \texttt{(click)="countClicks()"}
  \choice \texttt{(click)="count |= 1"}
  \choice \texttt{(click)="new Date()"}
\end{choices}
\begin{solution}
  Event bindings use parentheses: \lstinline[style=tsinline]{(click)="countClicks()"}. A typical handler lives in the component class:
\begin{lstlisting}[style=tsexam]
// app.component.ts
count = 0;
countClicks(): void {
  this.count += 1;
}
\end{lstlisting}
\begin{lstlisting}[style=htmlexam]
<button (click)="countClicks()">click me!</button>
<p>No of clicks: {{ count }}</p>
\end{lstlisting}
  Interpolation \lstinline[style=htmlexam]|{{ count }}| updates the view when \texttt{count} changes.
\end{solution}

\question[4]
Which syntax is \textbf{not} allowed in template \textbf{statements}?
\begin{choices}
  \choice A semicolon to chain two calls
  \choice A simple assignment with \texttt{=}
  \CorrectChoice Increment/decrement (\texttt{++}/\texttt{--}), operator assignment (\texttt{+=}), or the pipe operator \texttt{|}
  \choice Calling a component method
\end{choices}
\begin{solution}
  Template statements may chain with \texttt{;} or \texttt{,} and allow simple \texttt{=} assignment. They do \textbf{not} allow \texttt{++}, \texttt{--}, \texttt{+=}, \texttt{-=}, bitwise \texttt{\&}, or \texttt{|} (pipe). Keep handlers short---call methods on the component class.
\end{solution}

\question[4]
Which are realistic uses of template statements?
\begin{choices}
  \choice Submitting a form when the user clicks Submit
  \choice Disabling a button when input is empty
  \choice Incrementing a quantity with up/down controls
  \CorrectChoice All of the above
\end{choices}
\begin{solution}
  \textbf{All of the above.} Template statements wire UI events to component logic: submit handlers, conditional disabling, and increment/decrement actions---without putting complex logic inline in the HTML.
\end{solution}

\question[4]
Template \textbf{expressions} (used in interpolation and property binding) differ from template \textbf{statements} because expressions:
\begin{choices}
  \choice May use \texttt{++} and \texttt{new} freely
  \CorrectChoice Cannot use \texttt{new}, assignment, \texttt{++}/\texttt{--}, or chaining with \texttt{;}---they evaluate to a value in binding context
  \choice Run only on the server
  \choice Replace the need for a component class
\end{choices}
\begin{solution}
  \textbf{Expressions} (right-hand side of \lstinline[style=htmlexam]|{{ }}| or \lstinline[style=htmlexam]|[prop]="..."|) must stay side-effect free: no \texttt{new}, no \texttt{=}, no \texttt{++}/\texttt{--}, no \texttt{;}. \textbf{Statements} (event handlers) allow chaining and simple assignment but still forbid \texttt{++} and operator-assignment forms. Neither may reference \texttt{window} or \texttt{document} directly---use the component context.
\end{solution}

\newpage

\uplevel{\par\textbf{Part 12: One-way data binding}\par}

\question[4]
In Angular, \textbf{binding} is:
\begin{choices}
  \choice A Webpack plugin only
  \CorrectChoice Communication that synchronizes the template (view) and the component class (model), using change detection
  \choice A SQL connection string
  \choice The same as a CSS selector
\end{choices}
\begin{solution}
  Binding links template and model so the UI stays in sync with component data. Common \textbf{one-way} forms:
  \begin{itemize}
    \item \textbf{Text interpolation} --- model $\rightarrow$ template (\lstinline[style=htmlexam]|{{ }}|)
    \item \textbf{Property binding} --- model $\rightarrow$ element property (\lstinline[style=htmlexam]|[prop]="..."|)
    \item \textbf{Event binding} --- template $\rightarrow$ model (\lstinline[style=htmlexam]|(event)="..."|)
  \end{itemize}
  Angular also supports attribute, class, and style binding for accessibility and dynamic styling.
\end{solution}

\question[4]
\textbf{Text interpolation} transfers data:
\begin{choices}
  \choice From the template to the TypeScript class only
  \CorrectChoice From the TypeScript class (model) to the HTML template (view)
  \choice Only between two SQL tables
  \choice Without using the component class
\end{choices}
\begin{solution}
\begin{lstlisting}[style=htmlexam]
<p><b>{{ title }}</b></p>
\end{lstlisting}
\begin{lstlisting}[style=tsexam]
export class AppComponent {
  title = 'Template literal works!';
}
\end{lstlisting}
  The displayed text updates when \texttt{title} changes in the class. Interpolation syntax is double curly braces: \lstinline[style=htmlexam]|{{ }}|.
\end{solution}

\question[4]
\textbf{Property binding} is used to:
\begin{choices}
  \choice Listen for button clicks only
  \CorrectChoice Set or update an HTML element property from component data (for example an image \texttt{src})
  \choice Replace \texttt{NgModule} imports
  \choice Compile TypeScript inside the browser with no build step
\end{choices}
\begin{solution}
  Property binding flows model $\rightarrow$ template. Example: swap a logo by changing a variable in the class:
\begin{lstlisting}[style=htmlexam]
<img [src]="img_path">
\end{lstlisting}
\begin{lstlisting}[style=tsexam]
img_path = 'https://example.com/logo.png';
\end{lstlisting}
  Square brackets on the property name tell Angular to bind to the expression, not treat it as a plain string.
\end{solution}

\question[4]
\textbf{Event binding} is used to:
\begin{choices}
  \choice Display \lstinline[style=htmlexam]|{{ title }}| in the template
  \CorrectChoice Listen for user actions (clicks, keystrokes, etc.) and run component logic
  \choice Set the \texttt{src} attribute from the model
  \choice Import \texttt{FormsModule}
\end{choices}
\begin{solution}
  Event binding flows template $\rightarrow$ model:
\begin{lstlisting}[style=htmlexam]
<button (click)="onSave()">Save</button>
\end{lstlisting}
\begin{lstlisting}[style=tsexam]
onSave(): void {
  alert('Information Saved!');
}
\end{lstlisting}
  Parentheses on the event name wire the DOM event to a template statement or component method.
\end{solution}

\question[4]
Which binding type matches each task?
\begin{choices}
  \choice Clicks/submit $\rightarrow$ property binding; dynamic image $\rightarrow$ interpolation
  \CorrectChoice Clicks/submit $\rightarrow$ event binding; dynamic image $\rightarrow$ property binding; show a username from the class $\rightarrow$ text interpolation
  \choice All three tasks $\rightarrow$ event binding only
  \choice Dynamic image $\rightarrow$ text interpolation only
\end{choices}
\begin{solution}
  \textbf{Event binding} handles user input (quiz submit counts, save buttons). \textbf{Property binding} sets element properties such as \lstinline[style=htmlexam]|[src]|. \textbf{Interpolation} displays values like a user's name from the component field.
\end{solution}

\question[4]
Which syntax is used for \textbf{text interpolation}?
\begin{choices}
  \CorrectChoice \lstinline[style=htmlexam]|{{ }}|
  \choice \lstinline[style=htmlexam]|[]|
  \choice \lstinline[style=htmlexam]|{}|
  \choice \lstinline[style=htmlexam]|[()]|
\end{choices}
\begin{solution}
  \lstinline[style=htmlexam]|{{ expression }}| evaluates an expression from the component and inserts the result into the DOM. \lstinline[style=htmlexam]|[prop]="expr"| is property binding; \lstinline[style=htmlexam]|[(ngModel)]| is two-way binding.
\end{solution}

\newpage

\uplevel{\par\textbf{Part 13: Two-way data binding}\par}

\question[4]
\textbf{Two-way binding} in Angular means:
\begin{choices}
  \choice Data flows only from model to template, never back
  \CorrectChoice Data can flow from template to model and from model to template (view and class stay in sync)
  \choice Binding works only without a component class
  \choice The same as Webpack bundling
\end{choices}
\begin{solution}
  Two-way binding keeps a template field and a component property aligned. Angular's common syntax is \textbf{banana in a box}:
  \lstinline[style=htmlexam]|[(ngModel)]="propertyName"| ---combining property binding and event binding in one notation.
\end{solution}

\question[4]
To use \texttt{ngModel} for two-way binding, the module must:
\begin{choices}
  \choice Import \texttt{RouterModule} only
  \CorrectChoice Import \texttt{FormsModule} from \texttt{@angular/forms} in the \texttt{NgModule}
  \choice Delete \texttt{app.component.html}
  \choice Eject \texttt{webpack.config.js}
\end{choices}
\begin{solution}
\begin{lstlisting}[style=tsexam]
import { FormsModule } from '@angular/forms';

@NgModule({
  imports: [BrowserModule, FormsModule],
  // ...
})
export class AppModule { }
\end{lstlisting}
  Without \texttt{FormsModule}, \lstinline[style=htmlexam]|[(ngModel)]| is not available.
\end{solution}

\question[4]
In a two-way binding example, an input and a paragraph both show the same \texttt{value}, and a Clear button resets it. Which fragment is correct?
\begin{choices}
  \choice \lstinline[style=htmlexam]|<input [value]="value">| only
  \CorrectChoice \lstinline[style=htmlexam]|<input [(ngModel)]="value">| and \lstinline[style=htmlexam]|{{ value }}|
  \choice \lstinline[style=htmlexam]|{{ ngModel }}| only in the TypeScript file
  \choice \lstinline[style=htmlexam]|<button [click]="clearValue">| without parentheses
\end{choices}
\begin{solution}
\begin{lstlisting}[style=htmlexam]
<label>Enter Text:</label>
<input type="text" [(ngModel)]="value">
<p>{{ value }}</p>
<button (click)="clearValue()">Clear</button>
\end{lstlisting}
\begin{lstlisting}[style=tsexam]
value = '';
clearValue() {
  this.value = '';
}
\end{lstlisting}
  Typing updates \texttt{value} in the class; interpolation shows it; Clear sets \texttt{value} back to empty. Password-strength UIs use the same idea: rules check the model while the view shows feedback.
\end{solution}

\question[4]
Which methods can achieve \textbf{two-way binding}?
\begin{choices}
  \choice Using event binding and text interpolation/property binding together on the same field
  \choice Using \lstinline[style=htmlexam]|[(...)]| syntax (for example \lstinline[style=htmlexam]|[(ngModel)]|)
  \CorrectChoice Both A and B
  \choice None of the above
\end{choices}
\begin{solution}
  \textbf{Both.} You can combine one-way bindings manually (listen to \lstinline[style=htmlexam]|(input)| and bind \lstinline[style=htmlexam]|[value]|), or use \lstinline[style=htmlexam]|[(ngModel)]| as shorthand. \lstinline[style=htmlexam]|[(ngModel)]| is the usual approach for form fields.
\end{solution}

\question[4]
Summarize one-way versus two-way data flow:
\begin{choices}
  \choice Interpolation and property binding: template $\rightarrow$ model; events: model $\rightarrow$ template
  \CorrectChoice Interpolation and property binding: model $\rightarrow$ template; event binding: template $\rightarrow$ model; two-way: both directions
  \choice All bindings flow only template $\rightarrow$ model
  \choice Two-way binding removes the need for a component class
\end{choices}
\begin{solution}
  \textbf{One-way (model $\rightarrow$ view):} interpolation, property (and often attribute/class/style) binding.

  \textbf{One-way (view $\rightarrow$ model):} event binding.

  \textbf{Two-way:} \lstinline[style=htmlexam]|[(ngModel)]| (or equivalent property + event pair) keeps input and class property synchronized---useful for live form feedback such as password rules.
\end{solution}

\newpage

\uplevel{\par\textbf{Part 14: Structural directives}\par}

\question[4]
\textbf{Structural directives} in Angular are used to:
\begin{choices}
  \choice Change only the text color of an element without touching the DOM tree
  \CorrectChoice Add, remove, or manipulate DOM elements (not merely hide them with CSS in most cases)
  \choice Replace \texttt{NgModule} configuration
  \choice Compile TypeScript without a template
\end{choices}
\begin{solution}
  Structural directives reshape the DOM. They are written with a leading asterisk (for example \lstinline[style=htmlexam]|*ngIf|). Built-in examples:
  \begin{itemize}
    \item \lstinline[style=htmlexam]|*ngIf| --- conditionally include or exclude elements
    \item \lstinline[style=htmlexam]|*ngFor| --- repeat a template for each item in a list
    \item \lstinline[style=htmlexam]|*ngSwitch| / \lstinline[style=htmlexam]|*ngSwitchCase| / \lstinline[style=htmlexam]|*ngSwitchDefault| --- show one branch of markup
  \end{itemize}
  They can work with \lstinline[style=htmlexam]{<ng-template>} to hold template fragments.
\end{solution}

\question[4]
The \lstinline[style=htmlexam]|*ngIf| directive:
\begin{choices}
  \choice Only toggles \texttt{display:none} and always keeps the element in the DOM
  \CorrectChoice Adds or removes elements from the DOM based on a boolean expression
  \choice Iterates over an array in the template
  \choice Binds two-way form data like \lstinline[style=htmlexam]|[(ngModel)]|
\end{choices}
\begin{solution}
\begin{lstlisting}[style=htmlexam]
<p *ngIf="true">The expression is true; this paragraph is in the DOM.</p>
<p *ngIf="false">This paragraph is not in the DOM.</p>
\end{lstlisting}
  Use \lstinline[style=htmlexam]|*ngIf="condition; else elseBlock"| with \lstinline[style=htmlexam]{<ng-template \#elseBlock>} for an alternate block. Common uses: show content by user role, age gates, feature flags.
\end{solution}

\question[4]
The \lstinline[style=htmlexam]|*ngFor| directive is used to:
\begin{choices}
  \choice Switch among multiple template branches by matching a value
  \CorrectChoice Repeat a section of HTML once per item in an iterable list
  \choice Add CSS classes dynamically
  \choice Listen for click events only
\end{choices}
\begin{solution}
\begin{lstlisting}[style=tsexam]
customers: Customer[] = [
  { id: 234, name: 'John' },
  { id: 235, name: 'Adam' },
];
\end{lstlisting}
\begin{lstlisting}[style=htmlexam]
<tr *ngFor="let customer of customers">
  <td>{{ customer.id }}</td>
  <td>{{ customer.name }}</td>
</tr>
\end{lstlisting}
  After filtering with \lstinline[style=htmlexam]|ngSwitch|, a list of categories is often rendered with \lstinline[style=htmlexam]|*ngFor|.
\end{solution}

\question[4]
In \textbf{ngSwitch}, which pairing is correct?
\begin{choices}
  \choice \lstinline[style=htmlexam]|[ngSwitch]| on cases; \lstinline[style=htmlexam]|*ngSwitchCase| on the container
  \CorrectChoice \lstinline[style=htmlexam]|[ngSwitch]="expression"| on the container; \lstinline[style=htmlexam]|*ngSwitchCase| and \lstinline[style=htmlexam]|*ngSwitchDefault| on inner elements
  \choice \lstinline[style=htmlexam]|*ngSwitch| replaces \lstinline[style=htmlexam]|*ngFor| on table rows
  \choice \lstinline[style=htmlexam]|ngSwitchDefault| must appear before every case
\end{choices}
\begin{solution}
  \lstinline[style=htmlexam]|[ngSwitch]| is an attribute directive on the container that coordinates the case directives. \lstinline[style=htmlexam]|*ngSwitchCase| shows its element when its value matches; \lstinline[style=htmlexam]|*ngSwitchDefault| shows when no case matches---useful for filter menus and numeric input branches.
\begin{lstlisting}[style=htmlexam]
<div [ngSwitch]="num">
  <div *ngSwitchCase="'1'">You entered - One</div>
  <div *ngSwitchCase="'2'">You entered - Two</div>
  <div *ngSwitchDefault>...default...</div>
</div>
\end{lstlisting}
\end{solution}

\question[4]
Which of the following is \textbf{not} a structural directive?
\begin{choices}
  \choice \lstinline[style=htmlexam]|*ngIf|
  \choice \lstinline[style=htmlexam]|*ngFor|
  \choice \lstinline[style=htmlexam]|*ngSwitch|
  \CorrectChoice \lstinline[style=htmlexam]|[(ngModel)]|
\end{choices}
\begin{solution}
  \lstinline[style=htmlexam]|[(ngModel)]| is an \textbf{attribute}-style forms directive (two-way binding on inputs). Structural directives change which DOM nodes exist; \lstinline[style=htmlexam]|*ngIf| adds/removes by condition, \lstinline[style=htmlexam]|*ngFor| repeats by list.
\end{solution}

\question[4]
Which structural directive adds or removes DOM elements based on a condition?
\begin{choices}
  \CorrectChoice \lstinline[style=htmlexam]|*ngIf|
  \choice \lstinline[style=htmlexam]|*ngFor|
  \choice \lstinline[style=htmlexam]|*ngSwitch|
  \choice \lstinline[style=htmlexam]|[(ngModel)]|
\end{choices}
\begin{solution}
  \textbf{\lstinline[style=htmlexam]|*ngIf|} evaluates a boolean and includes or excludes the template content in the DOM. \lstinline[style=htmlexam]|*ngFor| loops; \lstinline[style=htmlexam]|ngSwitch| picks one branch among cases.
\end{solution}

\newpage

\uplevel{\par\textbf{Part 15: Attribute directives}\par}

\question[4]
\textbf{Attribute directives} are used to:
\begin{choices}
  \choice Remove entire sections of the DOM based on a list length
  \CorrectChoice Modify properties, appearance, or behaviour of existing DOM elements (classes, styles, etc.)
  \choice Replace the Angular compiler
  \choice Bootstrap the root \texttt{NgModule} only
\end{choices}
\begin{solution}
  Attribute directives change how an element looks or behaves without necessarily adding or removing nodes. Built-in examples include \lstinline[style=htmlexam]|[ngClass]|, \lstinline[style=htmlexam]|[ngStyle]|, and \lstinline[style=htmlexam]|[(ngModel)]|. Custom attribute directives share reusable UI logic (tabs, accordions, dropdowns) across components.
\end{solution}

\question[4]
The \lstinline[style=htmlexam]|[ngClass]| directive accepts:
\begin{choices}
  \choice Only a single hard-coded inline style string with no CSS classes
  \CorrectChoice A string, an array of class names, or an object map of class names to boolean expressions
  \choice Only numeric array indices from \lstinline[style=htmlexam]|*ngFor|
  \choice Only \lstinline[style=htmlexam]|(click)| event handlers
\end{choices}
\begin{solution}
\begin{lstlisting}[style=htmlexam]
<h3 [ngClass]="'red'">Need your attention</h3>
<div [ngClass]="['red', 'size20']">Red background, 20px italic text</div>
<div [ngClass]="{'red': false, 'size20': true}">Text with size20 only</div>
\end{lstlisting}
  Classes apply when the expression is truthy; object form toggles each class independently.
\end{solution}

\question[4]
The \lstinline[style=htmlexam]|[ngStyle]| directive is used to:
\begin{choices}
  \choice Repeat rows in a table from an array
  \CorrectChoice Apply inline styles dynamically from a component expression
  \choice Remove elements when a boolean is false
  \choice Match one branch in a switch statement
\end{choices}
\begin{solution}
\begin{lstlisting}[style=htmlexam]
<input type="text" [(ngModel)]="name">
<div [ngStyle]="{'background-color': name === 'Admin' ? 'green' : 'red'}">
</div>
\end{lstlisting}
  \lstinline[style=htmlexam]|[ngStyle]="objExp"| maps style property names to values computed in the template---useful for status colors and validation feedback.
\end{solution}

\question[4]
Which command creates a custom attribute directive named \texttt{highlight}?
\begin{choices}
  \choice \texttt{ng generate attribute highlight}
  \choice \texttt{ng create new attribute highlight}
  \CorrectChoice \texttt{ng generate directive highlight}
  \choice \texttt{ng new directive highlight}
\end{choices}
\begin{solution}
  The CLI adds \texttt{highlight.directive.ts} (and a spec file) and registers the class in the module. Example skeleton:
\begin{lstlisting}[style=tsexam]
@Directive({ selector: '[appHighlight]' })
export class HighlightDirective {
  constructor() { }
}
\end{lstlisting}
  Use the selector (for example \lstinline[style=htmlexam]|<p appHighlight>|) on elements in templates.
\end{solution}

\question[4]
Which of the following is an \textbf{attribute} directive?
\begin{choices}
  \choice \lstinline[style=htmlexam]|*ngIf|
  \CorrectChoice \lstinline[style=htmlexam]|[ngStyle]|
  \choice \lstinline[style=htmlexam]|*ngSwitch|
  \choice None of the above
\end{choices}
\begin{solution}
  \textbf{\lstinline[style=htmlexam]|[ngStyle]|} changes element styles. \lstinline[style=htmlexam]|*ngIf|, \lstinline[style=htmlexam]|*ngFor|, and \lstinline[style=htmlexam]|*ngSwitch| are \textbf{structural}---they alter the DOM layout, not only attributes on a fixed node.
\end{solution}

\question[4]
Attribute directives differ from structural directives mainly because attribute directives:
\begin{choices}
  \choice Always start with \texttt{*}
  \CorrectChoice Change properties or behaviour of elements that already exist, rather than adding or removing DOM nodes by default
  \choice Cannot be created with the CLI
  \choice Replace \lstinline[style=htmlexam]|{{ }}| interpolation
\end{choices}
\begin{solution}
  \textbf{Structural:} \lstinline[style=htmlexam]|*ngIf| / \lstinline[style=htmlexam]|*ngFor| / \lstinline[style=htmlexam]|ngSwitch| family reshapes the tree.

  \textbf{Attribute:} \lstinline[style=htmlexam]|[ngClass]|, \lstinline[style=htmlexam]|[ngStyle]|, \lstinline[style=htmlexam]|[(ngModel)]|, and custom directives such as \lstinline[style=htmlexam]|[appText]| adjust presentation or interaction on existing elements.
\end{solution}

\newpage

\uplevel{\par\textbf{Part 16: Pipes}\par}

\question[4]
In Angular templates, \textbf{pipes} are used to:
\begin{choices}
  \choice Compile TypeScript into JavaScript at build time
  \CorrectChoice Transform displayed data from one form into another (dates, currency, case, etc.)
  \choice Replace HTTP calls to the backend
  \choice Define root module metadata only
\end{choices}
\begin{solution}
  Backend data often arrives in a raw shape (ISO dates, plain numbers). Pipes format values \emph{in the view} without changing the underlying model. Examples: \lstinline[style=htmlexam]|date|, \lstinline[style=htmlexam]|currency|, \lstinline[style=htmlexam]|lowercase|, \lstinline[style=htmlexam]|percent|.
\end{solution}

\question[4]
Which is the correct \textbf{pipe syntax} in a template?
\begin{choices}
  \CorrectChoice \lstinline[style=htmlexam]!{{ text | pipeName }}!
  \choice \lstinline[style=htmlexam]![(text | pipeName)]!
  \choice \lstinline[style=htmlexam]|[[text : pipeName]]|
  \choice \lstinline[style=htmlexam]|((text : pipeName))|
\end{choices}
\begin{solution}
  Interpolation uses \lstinline[style=htmlexam]|{{ }}|; the vertical bar (\texttt{|}) separates the value from the pipe name. Optional parameters follow a colon: \lstinline[style=htmlexam]!{{ today | date:'MM-dd-yyyy' }}!.
\end{solution}

\question[4]
Angular pipes are classified into:
\begin{choices}
  \choice Structural pipes and attribute pipes only
  \CorrectChoice \textbf{Default} (built-in) pipes and \textbf{custom} pipes
  \choice One-way pipes and two-way pipes only
  \choice Module pipes and component pipes only
\end{choices}
\begin{solution}
  \textbf{Default pipes} ship with Angular (\lstinline[style=htmlexam]|date|, \lstinline[style=htmlexam]|currency|, \lstinline[style=htmlexam]|uppercase|, \lstinline[style=htmlexam]|percent|, and others). \textbf{Custom pipes} are classes you create when built-ins do not match your formatting rules.
\end{solution}

\question[4]
In a template, which expression formats a \texttt{Date} as \texttt{MM-dd-yyyy}?
\begin{choices}
  \choice \lstinline[style=htmlexam]!{{ today | lowercase }}!
  \CorrectChoice \lstinline[style=htmlexam]!{{ today | date:'MM-dd-yyyy' }}!
  \choice \lstinline[style=htmlexam]!{{ today | currency:'INR' }}!
  \choice \lstinline[style=htmlexam]!{{ today | titlecase }}!
\end{choices}
\begin{solution}
\begin{lstlisting}[style=tsexam]
today = new Date();
\end{lstlisting}
\begin{lstlisting}[style=htmlexam]
<p>Date: {{ today | date:'MM-dd-yyyy' }}</p>
\end{lstlisting}
  The \lstinline[style=htmlexam]|date| pipe turns an unreadable timestamp into a user-facing calendar string.
\end{solution}

\question[4]
Which template expression shows a numeric \texttt{price} with an explicit US dollar format?
\begin{choices}
  \CorrectChoice \lstinline[style=htmlexam]!{{ price | currency:'USD' }}!
  \choice \lstinline[style=htmlexam]!{{ price | uppercase }}!
  \choice \lstinline[style=htmlexam]!{{ price | date:'USD' }}!
  \choice \lstinline[style=htmlexam]!{{ price | titlecase }}!
\end{choices}
\begin{solution}
\begin{lstlisting}[style=htmlexam]
<p>Price: {{ price | currency:'INR' }}</p>
<p>USD: {{ price | currency:'USD' }}</p>
\end{lstlisting}
  \lstinline[style=htmlexam]|currency| adds the symbol and locale-aware grouping. \lstinline[style=htmlexam]|currency| without a code uses the default locale; \lstinline[style=htmlexam]|currency:'USD'| forces US dollars.
\end{solution}

\question[4]
Which built-in pipes change \textbf{text case} in the template?
\begin{choices}
  \choice \lstinline[style=htmlexam]|date| and \lstinline[style=htmlexam]|currency| only
  \CorrectChoice \lstinline[style=htmlexam]|lowercase|, \lstinline[style=htmlexam]|uppercase|, and \lstinline[style=htmlexam]|titlecase|
  \choice \lstinline[style=htmlexam]|*ngIf| and \lstinline[style=htmlexam]|*ngFor| only
  \choice \lstinline[style=htmlexam]|json| and \lstinline[style=htmlexam]|async| only
\end{choices}
\begin{solution}
\begin{lstlisting}[style=htmlexam]
<p>{{ title | lowercase }}</p>
<p>{{ title | uppercase }}</p>
<p>{{ title | titlecase }}</p>
\end{lstlisting}
  Use these when labels, headings, or user-visible strings must follow a consistent casing rule.
\end{solution}

\uplevel{\par\textbf{Part 17: Custom pipes}\par}

\question[4]
A \textbf{custom pipe} is most appropriate when:
\begin{choices}
  \choice The built-in \lstinline[style=htmlexam]|date| pipe already covers every display rule you need
  \CorrectChoice Default pipes cannot produce the required format (for example initials from an email address)
  \choice You need to bootstrap \texttt{AppModule}
  \choice You want to remove DOM nodes conditionally
\end{choices}
\begin{solution}
  Custom pipes implement \lstinline[style=tsinline]|PipeTransform| for domain-specific formatting---masked names, internal IDs, compact avatars from email, and similar rules that are not covered by \lstinline[style=htmlexam]|date|/\lstinline[style=htmlexam]|currency|/\lstinline[style=htmlexam]|case| pipes.
\end{solution}

\question[4]
Which CLI command creates a pipe named \texttt{user}?
\begin{choices}
  \choice \texttt{ng create pipe user}
  \choice \texttt{generate pipe user}
  \choice \texttt{pipe user}
  \CorrectChoice \texttt{ng generate pipe user}
\end{choices}
\begin{solution}
  The CLI generates \texttt{user.pipe.ts} (and a spec file) and registers the pipe in the module \texttt{declarations} array. Template usage: \lstinline[style=htmlexam]!{{ userName | user }}!.
\end{solution}

\question[4]
A custom pipe class must:
\begin{choices}
  \choice Extend \texttt{Component} and declare a \texttt{templateUrl}
  \CorrectChoice Be decorated with \lstinline[style=tsinline]|@Pipe| and implement \lstinline[style=tsinline]|PipeTransform|
  \choice Replace \lstinline[style=htmlexam]|[(ngModel)]| on all inputs
  \choice Be listed only in \texttt{bootstrap}, never in \texttt{declarations}
\end{choices}
\begin{solution}
\begin{lstlisting}[style=tsexam]
@Pipe({ name: 'user' })
export class UserPipe implements PipeTransform {
  transform(value: string, ...args: any[]): string {
    const values = value.split('.');
    return values[0][0] + values[1][0];
  }
}
\end{lstlisting}
  Angular calls \lstinline[style=tsinline]|transform()| whenever the pipe appears in a template binding.
\end{solution}

\question[4]
Given \lstinline[style=tsinline]|userName = 'taylor.swift@xyz.com'|, what does this template produce after both pipes run?
\begin{lstlisting}[style=htmlexam]
<p>{{ userName | user | uppercase }}</p>
\end{lstlisting}
\begin{choices}
  \choice The full email in lowercase
  \choice \texttt{taylor.swift} in title case
  \CorrectChoice \texttt{TS} (initials from the local part, then uppercased)
  \choice An empty string always
\end{choices}
\begin{solution}
  The custom \lstinline[style=htmlexam]|user| pipe splits on \texttt{.}, takes the first character of each segment (\texttt{t}, \texttt{s}), then \lstinline[style=htmlexam]|uppercase| yields \texttt{TS}. Pipes chain left to right: each pipe's output becomes the next pipe's input.
\end{solution}

\question[4]
To use a custom pipe in a component template, you typically:
\begin{choices}
  \choice Import it only in \texttt{main.ts}, not in the module
  \CorrectChoice Declare it in the module's \texttt{declarations} (or import a module that exports it), then use \lstinline[style=htmlexam]!{{ value | pipeName }}!
  \choice Add it to the \texttt{bootstrap} array instead of \texttt{declarations}
  \choice Prefix the pipe name with \texttt{*} in HTML
\end{choices}
\begin{solution}
  Like components and directives, pipes belong to an \texttt{NgModule}. Export the pipe from a shared module when other feature modules need the same formatter.
\end{solution}

\newpage

\uplevel{\par\textbf{Part 18: \texttt{@NgModule} decorator}\par}

\question[4]
An Angular \textbf{NgModule} is:
\begin{choices}
  \choice A CSS file that styles the root component only
  \CorrectChoice A TypeScript class marked with \lstinline[style=tsinline]|@NgModule| that groups components, directives, and pipes and configures compilation
  \choice The same as a single \texttt{*ngFor} loop
  \choice A Webpack loader configuration file
\end{choices}
\begin{solution}
  Every app has at least one module (usually \texttt{AppModule}). The decorator metadata tells Angular which declarations belong to the module, which other modules to import, which services to provide, and which component to bootstrap at startup.
\end{solution}

\question[4]
Which decorator supplies metadata for the \textbf{root module}?
\begin{choices}
  \CorrectChoice \lstinline[style=tsinline]|@NgModule|
  \choice \lstinline[style=tsinline]|@Module|
  \choice \lstinline[style=tsinline]|@Component| only
  \choice \lstinline[style=tsinline]|@Injectable| only
\end{choices}
\begin{solution}
  \lstinline[style=tsinline]|@NgModule| is imported from \lstinline[style=tsinline]|@angular/core|. \lstinline[style=tsinline]|@Component| describes one view class; \lstinline[style=tsinline]|@Injectable| marks a service for dependency injection.
\end{solution}

\question[4]
Which of the following is a valid \lstinline[style=tsinline]|@NgModule| metadata property?
\begin{choices}
  \choice \texttt{definitions}
  \choice \texttt{privileges}
  \CorrectChoice \texttt{declarations}
  \choice \texttt{services}
\end{choices}
\begin{solution}
  Common metadata keys include \texttt{declarations}, \texttt{imports}, \texttt{exports}, \texttt{providers}, and \texttt{bootstrap}. There is no \texttt{definitions}, \texttt{privileges}, or \texttt{services} key---service classes are registered under \texttt{providers}.
\end{solution}

\question[4]
In a typical \texttt{AppModule}, the \texttt{declarations} array lists:
\begin{choices}
  \choice Only third-party npm packages
  \CorrectChoice Components, directives, and pipes that belong to this module
  \choice The entry point file \texttt{main.ts}
  \choice Global CSS variables only
\end{choices}
\begin{solution}
  \texttt{declarations} registers classes owned by the module. \texttt{imports} brings in other modules (for example \texttt{BrowserModule}). \texttt{exports} makes selected declarations available to importing modules.
\end{solution}

\question[4]
The \texttt{bootstrap} property of \texttt{@NgModule} identifies:
\begin{choices}
  \choice Every route in the router configuration
  \CorrectChoice The root component Angular loads when the application starts
  \choice All HTTP interceptors in the project
  \choice The contents of \texttt{angular.json} only
\end{choices}
\begin{solution}
  CLI-generated apps place \texttt{AppComponent} in both \texttt{declarations} and \texttt{bootstrap}. \texttt{main.ts} calls \lstinline[style=tsinline]|platformBrowserDynamic().bootstrapModule(AppModule)|, which then creates the bootstrap component and inserts its template into \lstinline[style=htmlexam]|<app-root>| in \texttt{index.html}.
\end{solution}

\question[4]
Which sequence best describes the default Angular \textbf{startup flow}?
\begin{choices}
  \choice \texttt{index.html} $\rightarrow$ \texttt{angular.json} $\rightarrow$ \texttt{AppComponent} only
  \CorrectChoice \texttt{angular.json} (build entry) $\rightarrow$ \texttt{main.ts} $\rightarrow$ \texttt{AppModule} $\rightarrow$ \texttt{AppComponent} $\rightarrow$ \texttt{index.html} / \lstinline[style=htmlexam]|<app-root>|
  \choice \texttt{AppComponent} $\rightarrow$ \texttt{main.ts} with no module
  \choice \texttt{package.json} $\rightarrow$ pipes $\rightarrow$ \texttt{AppModule}
\end{choices}
\begin{solution}
  \texttt{angular.json} points the builder at \texttt{main.ts}. \texttt{main.ts} bootstraps \texttt{AppModule}; the module's \texttt{bootstrap} array starts \texttt{AppComponent}, whose selector hosts the root template inside \lstinline[style=htmlexam]|<app-root>|.
\end{solution}

\newpage

\uplevel{\par\textbf{Part 19: Event emitters}\par}

\question[4]
In a parent--child component relationship, data usually flows \textbf{down} with \lstinline[style=tsinline]|@Input| and flows \textbf{up} with:
\begin{choices}
  \choice \lstinline[style=htmlexam]|*ngFor| only
  \CorrectChoice \lstinline[style=tsinline]|@Output| and an \lstinline[style=tsinline]|EventEmitter|
  \choice \lstinline[style=htmlexam]|[(ngModel)]| on the parent selector only
  \choice \texttt{bootstrap} in \texttt{AppModule}
\end{choices}
\begin{solution}
  \textbf{Parent $\rightarrow$ child:} bind a property with \lstinline[style=htmlexam]|[inputProp]="parentValue"| on the child selector (backed by \lstinline[style=tsinline]|@Input()| in the child class).

  \textbf{Child $\rightarrow$ parent:} the child raises a custom event; the parent listens in its template with event binding.
\end{solution}

\question[4]
The \lstinline[style=tsinline]|@Output| decorator marks a component property as:
\begin{choices}
  \choice A route parameter read from the URL
  \CorrectChoice A custom event channel the parent can subscribe to in its template
  \choice A CSS class applied with \lstinline[style=htmlexam]|[ngClass]|
  \choice A pipe registered in \texttt{declarations}
\end{choices}
\begin{solution}
  \lstinline[style=tsinline]|@Output| exposes an \lstinline[style=tsinline]|EventEmitter| instance. The parent binds to the \emph{output property name} (often aliased with \lstinline[style=tsinline]|@Output('alias')|) using the same \lstinline[style=htmlexam]|(eventName)="handler()"| syntax as native DOM events.
\end{solution}

\question[4]
Which child component code correctly declares an output that sends a \texttt{string} to the parent?
\begin{choices}
  \choice \lstinline[style=tsinline]|@Input() saved = new EventEmitter<string>();|
  \CorrectChoice \lstinline[style=tsinline]|@Output() saved = new EventEmitter<string>();|
  \choice \lstinline[style=tsinline]|@Output saved: string;|
  \choice \lstinline[style=tsinline]|@Injectable() saved = EventEmitter;|
\end{choices}
\begin{solution}
\begin{lstlisting}[style=tsexam]
import { Component, Output, EventEmitter } from '@angular/core';

@Component({ selector: 'app-save-button', template: `...` })
export class SaveButtonComponent {
  @Output() saved = new EventEmitter<string>();

  onClick() {
    this.saved.emit('draft-42');
  }
}
\end{lstlisting}
  Import \lstinline[style=tsinline]|EventEmitter| from \lstinline[style=tsinline]|@angular/core| and call \lstinline[style=tsinline]|emit()| with the payload.
\end{solution}

\question[4]
Given a child selector \lstinline[style=htmlexam]|<app-item (picked)="onPicked($event)">|, the parent receives the payload in:
\begin{choices}
  \choice \lstinline[style=htmlexam]|{{ picked }}| automatically, with no method
  \CorrectChoice The \lstinline[style=tsinline]|$event| argument passed to \lstinline[style=tsinline]|onPicked|
  \choice The \texttt{providers} array of \texttt{AppModule}
  \choice \lstinline[style=htmlexam]|[ngModel]| on the child only
\end{choices}
\begin{solution}
  Custom outputs behave like DOM events in templates: \lstinline[style=htmlexam]|(picked)="onPicked($event)"| forwards whatever the child passed to \lstinline[style=tsinline]|emit()|. Example handler:
\begin{lstlisting}[style=tsexam]
onPicked(id: string) {
  this.selectedId = id;
}
\end{lstlisting}
\end{solution}

\question[4]
To notify the parent that a form was submitted, the child should:
\begin{choices}
  \choice Mutate a parent field directly with \lstinline[style=tsinline]|this.parent.title = 'done'|
  \CorrectChoice Call \lstinline[style=tsinline]|this.formSubmit.emit(formValue)| on its \lstinline[style=tsinline]|EventEmitter|
  \choice Call \lstinline[style=tsinline]|bootstrapModule| again
  \choice Add the form to \texttt{declarations} twice
\end{choices}
\begin{solution}
  Children should not reach into parent instances. Emitting an event keeps components loosely coupled: the parent decides whether to save, navigate, or show a message.
\end{solution}

\question[4]
Which template wiring is valid for a child \lstinline[style=htmlexam]|<app-counter>| that exposes \lstinline[style=tsinline]|@Output() countChange|?
\begin{choices}
  \choice \lstinline[style=htmlexam]|<app-counter [countChange]="total">|
  \CorrectChoice \lstinline[style=htmlexam]|<app-counter (countChange)="updateTotal($event)">|
  \choice \lstinline[style=htmlexam]|<app-counter [(countChange)]="total">| without \lstinline[style=tsinline]|@Input|
  \choice \lstinline[style=htmlexam]|<app-counter *countChange="total">|
\end{choices}
\begin{solution}
  Square brackets \lstinline[style=htmlexam]|[]| set \textbf{inputs}; parentheses \lstinline[style=htmlexam]|()| listen for \textbf{outputs}. Two-way \lstinline[style=htmlexam]|[( )]| on a custom property requires both \lstinline[style=tsinline]|@Input| and \lstinline[style=tsinline]|@Output| with a coordinated naming pattern.
\end{solution}

\uplevel{\par\textbf{Part 20: Component styles}\par}

\question[4]
By default, styles listed in \texttt{styleUrls} for a component are:
\begin{choices}
  \choice Applied globally to every component with no scoping
  \CorrectChoice Scoped to that component's template (emulated encapsulation by default)
  \choice Ignored unless placed in \texttt{angular.json} only
  \choice Compiled into SQL queries
\end{choices}
\begin{solution}
  Angular adds unique attributes to elements and adjusts selectors so rules from \texttt{hero.component.css} do not leak into sibling components. Global rules belong in \texttt{styles.css} or the \texttt{styles} array in \texttt{angular.json}.
\end{solution}

\question[4]
Which \texttt{@Component} metadata attaches external CSS files to a component?
\begin{choices}
  \choice \texttt{providers}
  \choice \texttt{encapsulation: 'css'}
  \CorrectChoice \texttt{styleUrls: ['./app.component.css']}
  \choice \texttt{bootstrap}
\end{choices}
\begin{solution}
  Use \texttt{styleUrls} for one or more file paths, or inline \texttt{styles: [\ldots]} in the decorator for small rules. Both pair with the component's template the same way \texttt{templateUrl} does for HTML.
\end{solution}

\question[4]
Setting \lstinline[style=tsinline]|encapsulation: ViewEncapsulation.None| means:
\begin{choices}
  \choice Styles apply only inside Shadow DOM and never reach the page
  \CorrectChoice Component styles become global (no emulated scoping attributes)
  \choice The component cannot use CSS classes
  \choice \lstinline[style=htmlexam]|*ngIf| is disabled
\end{choices}
\begin{solution}
  \textbf{Emulated} (default) scopes styles to the template. \textbf{None} turns off scoping---useful for third-party widgets or legacy CSS, but risks clashes with other components. \textbf{ShadowDom} uses native shadow roots where supported.
\end{solution}

\question[4]
The \lstinline[style=htmlexam]|:host| selector in a component stylesheet targets:
\begin{choices}
  \choice Every \texttt{<p>} tag in the entire application
  \CorrectChoice The host element of the component (for example \lstinline[style=htmlexam]|<app-hero>|)
  \choice Only elements inside \lstinline[style=htmlexam]|*ngFor| loops
  \choice The \texttt{RouterOutlet} exclusively
\end{choices}
\begin{solution}
\begin{lstlisting}[style=htmlexam]
:host {
  display: block;
  border: 1px solid #ccc;
}
:host(.active) {
  border-color: green;
}
\end{lstlisting}
  \lstinline[style=htmlexam]|:host-context(.theme-dark)| applies when an ancestor matches (useful for theming).
\end{solution}

\question[4]
Where should application-wide styles (site fonts, reset rules) normally live?
\begin{choices}
  \choice Only inside each component's \texttt{styleUrls}, duplicated per file
  \CorrectChoice In global styles such as \texttt{src/styles.css} (listed under \texttt{styles} in \texttt{angular.json})
  \choice In the \texttt{bootstrap} array of \texttt{@NgModule}
  \choice In \lstinline[style=tsinline]|EventEmitter| definitions
\end{choices}
\begin{solution}
  Component CSS handles layout and look for one view. Global CSS sets baselines shared by the whole app. Keep globals small to avoid overriding encapsulated component rules unintentionally.
\end{solution}

\question[4]
Component styles with default encapsulation differ from global styles mainly because encapsulated rules:
\begin{choices}
  \choice Run on the server only during \texttt{ng build}
  \CorrectChoice Match only elements from that component's template (plus \lstinline[style=htmlexam]|:host| / \lstinline[style=htmlexam]|:host-context|)
  \choice Require \lstinline[style=tsinline]|@Injectable|
  \choice Replace TypeScript compilation
\end{choices}
\begin{solution}
  Encapsulation prevents a button style in \texttt{dashboard.component.css} from restyling unrelated buttons in \texttt{settings.component.html}. When you truly need to style a child from a parent, prefer documented patterns (\lstinline[style=htmlexam]|:host ::ng-deep| is deprecated---restructure styles or use \texttt{ViewEncapsulation.None} sparingly).
\end{solution}

\newpage

\uplevel{\par\textbf{Part 21: Route guards}\par}

\question[4]
A \textbf{route guard} in Angular is used to:
\begin{choices}
  \choice Compile templates into JavaScript only
  \CorrectChoice Allow, block, or redirect navigation to a route (for example login checks)
  \choice Replace \lstinline[style=htmlexam]|{{ }}| interpolation
  \choice Format dates with the \lstinline[style=htmlexam]|date| pipe
\end{choices}
\begin{solution}
  Guards run during routing. Common interfaces include \lstinline[style=tsinline]|CanActivate| (enter a route), \lstinline[style=tsinline]|CanDeactivate| (leave a route), \lstinline[style=tsinline]|CanMatch| (whether a route config applies), and \lstinline[style=tsinline]|Resolve| (fetch data before activation).
\end{solution}

\question[4]
The \lstinline[style=tsinline]|CanActivate| guard runs:
\begin{choices}
  \choice After the user closes the browser tab only
  \CorrectChoice Before a route is activated, to decide if navigation may proceed
  \choice Only when a component calls \lstinline[style=tsinline]|emit()|
  \choice During CSS encapsulation
\end{choices}
\begin{solution}
  If \lstinline[style=tsinline]|canActivate| returns \texttt{false} (or a \lstinline[style=tsinline]|UrlTree| redirect), Angular cancels or alters navigation. Typical checks: authentication, roles, feature flags.
\end{solution}

\question[4]
Which route configuration applies an \texttt{AuthGuard} before loading \texttt{AdminComponent}?
\begin{choices}
  \choice \lstinline[style=tsinline]|{ path: 'admin', providers: [AuthGuard] }|
  \CorrectChoice \lstinline[style=tsinline]|{ path: 'admin', component: AdminComponent, canActivate: [AuthGuard] }|
  \choice \lstinline[style=tsinline]|{ path: 'admin', bootstrap: [AuthGuard] }|
  \choice \lstinline[style=tsinline]|{ path: 'admin', declarations: [AuthGuard] }|
\end{choices}
\begin{solution}
  Register guards on the route (or parent route). The guard class is usually \lstinline[style=tsinline]|@Injectable({ providedIn: 'root' })| and listed in \lstinline[style=tsinline]|canActivate| (or \lstinline[style=tsinline]|canActivateChild| for child routes).
\end{solution}

\question[4]
An \texttt{AuthGuard} implementation should:
\begin{choices}
  \choice Extend \lstinline[style=tsinline]|Component| and declare a \texttt{templateUrl}
  \CorrectChoice Implement \lstinline[style=tsinline]|CanActivate| and return a \texttt{boolean}, \lstinline[style=tsinline]|UrlTree|, or an \texttt{Observable}/\texttt{Promise} of those types
  \choice Implement \lstinline[style=tsinline]|PipeTransform| only
  \choice Be placed in the \texttt{declarations} array as a component
\end{choices}
\begin{solution}
\begin{lstlisting}[style=tsexam]
@Injectable({ providedIn: 'root' })
export class AuthGuard implements CanActivate {
  constructor(private auth: AuthService, private router: Router) {}

  canActivate(): boolean | UrlTree {
    if (this.auth.isLoggedIn()) {
      return true;
    }
    return this.router.createUrlTree(['/login']);
  }
}
\end{lstlisting}
  Returning a \lstinline[style=tsinline]|UrlTree| sends the user to login without manually calling \lstinline[style=tsinline]|navigate| inside the guard.
\end{solution}

\question[4]
A \lstinline[style=tsinline]|CanDeactivate| guard is commonly used to:
\begin{choices}
  \choice Add currency symbols to prices
  \CorrectChoice Ask whether the user may leave a route with unsaved form changes
  \choice Bootstrap \texttt{AppModule}
  \choice Scope CSS with \lstinline[style=htmlexam]|:host|
\end{choices}
\begin{solution}
  Example: a profile editor calls \lstinline[style=tsinline]|canDeactivate()| on the component or a guard service. If the form is dirty, show a confirmation dialog; return \texttt{false} to stay on the page.
\end{solution}

\question[4]
If \lstinline[style=tsinline]|canActivate| returns \texttt{false} and no \lstinline[style=tsinline]|UrlTree| is returned, navigation:
\begin{choices}
  \CorrectChoice Is cancelled (the current route stays active)
  \choice Always redirects to \texttt{/admin} automatically
  \choice Deletes all \texttt{providers} from the module
  \choice Recompiles every component with \texttt{ViewEncapsulation.None}
\end{choices}
\begin{solution}
  Return \texttt{true} to allow activation. Return \texttt{false} to block. Return a \lstinline[style=tsinline]|UrlTree| (for example \lstinline[style=tsinline]|router.createUrlTree(['/login'])|) to redirect. Async guards may return \lstinline[style=tsinline]|Observable<boolean | UrlTree>| when waiting on an HTTP auth check.
\end{solution}

\uplevel{\par\textbf{Part 22: Dependency injection}\par}

\question[4]
\textbf{Dependency injection} (DI) in Angular means:
\begin{choices}
  \choice Hard-coding \texttt{new Service()} in every template
  \CorrectChoice The framework supplies required classes (services) to constructors or \lstinline[style=tsinline]|inject()| call sites from injectors
  \choice Importing CSS into \texttt{styleUrls}
  \choice Emitting DOM events with \lstinline[style=tsinline]|EventEmitter|
\end{choices}
\begin{solution}
  Components and services declare \emph{dependencies} (for example \texttt{AuthService}, \texttt{HttpClient}). Angular's injector creates or reuses instances and passes them in---improving testability (mock services) and keeping construction logic out of templates.
\end{solution}

\question[4]
Which decorator marks a class as a service that can be injected?
\begin{choices}
  \choice \lstinline[style=tsinline]|@Component| only
  \CorrectChoice \lstinline[style=tsinline]|@Injectable|
  \choice \lstinline[style=tsinline]|@Pipe| only
  \choice \lstinline[style=tsinline]|@NgModule| on every service class
\end{choices}
\begin{solution}
  Services are plain TypeScript classes plus \lstinline[style=tsinline]|@Injectable|. Without it, Angular may not be able to generate the metadata needed for constructor injection in some build modes.
\end{solution}

\question[4]
\lstinline[style=tsinline]|@Injectable({ providedIn: 'root' })| registers the service:
\begin{choices}
  \choice Separately in every component's \texttt{styleUrls}
  \CorrectChoice With the root injector as a singleton shared across the application
  \choice Only inside \lstinline[style=htmlexam]|*ngFor| templates
  \choice As a route guard automatically
\end{choices}
\begin{solution}
  \texttt{providedIn: 'root'} (tree-shakable) is the common pattern for app-wide services. Alternatives: list the service in a module's \texttt{providers} array, or \lstinline[style=tsinline]|providedIn: 'platform'| for multi-app scenarios.
\end{solution}

\question[4]
Which constructor correctly requests an injected \texttt{UserService}?
\begin{choices}
  \choice \lstinline[style=tsinline]|constructor() { this.user = new UserService(); }|
  \CorrectChoice \lstinline[style=tsinline]|constructor(private user: UserService) {}|
  \choice \lstinline[style=tsinline]|constructor(user: UserService = null) {}| with no decorator
  \choice \lstinline[style=tsinline]|constructor(@Output() user: UserService) {}|
\end{choices}
\begin{solution}
  Parameter properties (\lstinline[style=tsinline]|private user: UserService|) tell Angular to inject \texttt{UserService} when creating the class. The same service instance is reused according to where it was provided (root vs component vs module).
\end{solution}

\question[4]
Registering a service in an \texttt{NgModule}'s \texttt{providers} array:
\begin{choices}
  \choice Automatically registers the service as a route guard
  \CorrectChoice Supplies the service to the module injector (and typically its components), which may differ from \texttt{providedIn: 'root'} scope
  \choice Replaces \lstinline[style=tsinline]|EventEmitter| on child components
  \choice Is required for every pipe name
\end{choices}
\begin{solution}
  \textbf{Root} (\lstinline[style=tsinline]|providedIn: 'root'|): one instance for the whole app. \textbf{Module/component providers}: new instance per injector that declares them---useful for feature-specific state. Avoid providing the same service in multiple places unless you intend separate instances.
\end{solution}

\question[4]
In modern Angular, \lstinline[style=tsinline]|inject(DataService)| is used to:
\begin{choices}
  \choice Replace all HTML templates with JSON
  \CorrectChoice Obtain a service reference inside a factory, functional guard, or \texttt{runInInjectionContext} without a constructor parameter
  \choice Emit events to parent components
  \choice Set \lstinline[style=htmlexam]|ViewEncapsulation.None|
\end{choices}
\begin{solution}
  Constructor injection remains the default in components. The \lstinline[style=tsinline]|inject()| function from \lstinline[style=tsinline]|@angular/core| fits functional route guards, interceptors, and standalone bootstrapping where a class constructor is awkward. It must run in an injection context (for example during construction or an \texttt{EnvironmentInjector} callback).
\end{solution}

\newpage

\uplevel{\par\textbf{Part 23: \texttt{HttpClient}}\par}

\question[4]
Angular's \texttt{HttpClient} is used to:
\begin{choices}
  \choice Compile SCSS into component CSS only
  \CorrectChoice Send HTTP requests to REST APIs and receive responses (often JSON) as reactive streams
  \choice Replace \lstinline[style=tsinline]|@NgModule| metadata
  \choice Scope styles with \lstinline[style=htmlexam]|:host|
\end{choices}
\begin{solution}
  \texttt{HttpClient} (from \lstinline[style=tsinline]|@angular/common/http|) wraps the browser \texttt{XMLHttpRequest}/\texttt{fetch} APIs. Typical uses: load lists, save forms, authenticate, upload files. Responses are usually parsed as JavaScript objects when the server returns JSON.
\end{solution}

\question[4]
To use \texttt{HttpClient} in a classic \texttt{NgModule} app, you must:
\begin{choices}
  \choice Add \texttt{HttpClient} to the \texttt{bootstrap} array only
  \CorrectChoice Import \lstinline[style=tsinline]|HttpClientModule| (or provide \lstinline[style=tsinline]|provideHttpClient()| in standalone/bootstrap setups) and inject \lstinline[style=tsinline]|HttpClient|
  \choice Declare \texttt{HttpClient} in \texttt{declarations}
  \choice Call \lstinline[style=tsinline]|emit()| on every GET request
\end{choices}
\begin{solution}
\begin{lstlisting}[style=tsexam]
import { HttpClientModule } from '@angular/common/http';

@NgModule({
  imports: [BrowserModule, HttpClientModule],
})
export class AppModule {}
\end{lstlisting}
  Inject \lstinline[style=tsinline]|HttpClient| into services or components: \lstinline[style=tsinline]|constructor(private http: HttpClient) {}|.
\end{solution}

\question[4]
The return type of \lstinline[style=tsinline]|this.http.get<Hero[]>('/api/heroes')| is:
\begin{choices}
  \choice A \texttt{Promise<Hero[]>} only
  \choice A plain \texttt{Hero[]} immediately
  \CorrectChoice An \lstinline[style=tsinline]|Observable<Hero[]>| that emits when the response arrives
  \choice A \lstinline[style=tsinline]|EventEmitter<Hero[]>|
\end{choices}
\begin{solution}
  \texttt{HttpClient} methods return \textbf{Observables} (lazy---nothing is sent until you \texttt{subscribe} or use the \lstinline[style=htmlexam]|async| pipe). You can transform streams with RxJS operators (\lstinline[style=tsinline]|map|, \lstinline[style=tsinline]|catchError|, \lstinline[style=tsinline]|retry|).
\end{solution}

\question[4]
Which code correctly loads heroes into a component field?
\begin{choices}
  \choice \lstinline[style=tsinline]|this.heroes = this.http.get('/api/heroes');| (no subscribe)
  \CorrectChoice \lstinline[style=tsinline]|this.http.get<Hero[]>('/api/heroes').subscribe(data => this.heroes = data);|
  \choice \lstinline[style=tsinline]|@Output() heroes = this.http.get(...);|
  \choice \lstinline[style=htmlexam]|{{ http.get('/api/heroes') }}| in the template without a pipe
\end{choices}
\begin{solution}
  Subscribe (or bind with \lstinline[style=htmlexam]|heroes$ | async| in the template) to trigger the request and receive data. Unsubscribing on destroy avoids leaks when not using \lstinline[style=htmlexam]|async|.
\end{solution}

\question[4]
To send JSON in a POST body with \texttt{HttpClient}, you typically:
\begin{choices}
  \choice Call \lstinline[style=tsinline]|http.post(url)| with no body and put JSON in \texttt{styleUrls}
  \CorrectChoice Call \lstinline[style=tsinline]|http.post(url, body, { headers })| where \texttt{body} is a JavaScript object
  \choice Use \lstinline[style=htmlexam]|*ngFor| on the URL string
  \choice Register the POST in \texttt{declarations}
\end{choices}
\begin{solution}
\begin{lstlisting}[style=tsexam]
this.http.post<Hero>('/api/heroes', hero, {
  headers: new HttpHeaders({ 'Content-Type': 'application/json' }),
}).subscribe(saved => console.log(saved));
\end{lstlisting}
  \texttt{HttpClient} serializes plain objects to JSON by default for write operations.
\end{solution}

\question[4]
When an HTTP request fails (network error or \texttt{4xx}/\texttt{5xx}), handling it in the service layer often uses:
\begin{choices}
  \choice \lstinline[style=tsinline]|ViewEncapsulation.None|
  \CorrectChoice RxJS \lstinline[style=tsinline]|catchError| (and sometimes \lstinline[style=tsinline]|retry|) on the Observable pipeline
  \choice \lstinline[style=htmlexam]|[(ngModel)]| on the error object
  \choice \texttt{Karma} \lstinline[style=tsinline]|describe| blocks only
\end{choices}
\begin{solution}
  \lstinline[style=tsinline]|catchError| maps failures to a fallback value or rethrows a friendly error. Components can show a message while the service keeps HTTP details centralized---easier to unit test with \texttt{HttpTestingController}.
\end{solution}

\uplevel{\par\textbf{Part 24: Pub--sub design pattern}\par}

\question[4]
The \textbf{publish--subscribe (pub--sub)} pattern means:
\begin{choices}
  \choice Every component must call every other component's methods directly
  \CorrectChoice Publishers emit messages without knowing subscribers; subscribers react without knowing publishers
  \choice CSS is applied only with \lstinline[style=htmlexam]|:host|
  \choice Routes are defined only in \texttt{index.html}
\end{choices}
\begin{solution}
  Pub--sub \textbf{decouples} producers and consumers. In Angular, a shared injectable service often holds a \lstinline[style=tsinline]|Subject| or \lstinline[style=tsinline]|BehaviorSubject|; components \texttt{publish} with \lstinline[style=tsinline]|next()| and \texttt{subscribe} to updates.
\end{solution}

\question[4]
How does pub--sub differ from a child \lstinline[style=tsinline]|@Output()| \lstinline[style=tsinline]|EventEmitter|?
\begin{choices}
  \CorrectChoice \lstinline[style=tsinline]|EventEmitter| is for parent--child communication; pub--sub usually broadcasts through a shared service to any interested listener
  \choice They are identical and interchangeable in all cases
  \choice \lstinline[style=tsinline]|EventEmitter| only works with \texttt{HttpClient}
  \choice Pub--sub requires \lstinline[style=htmlexam]|*ngIf| on the root element
\end{choices}
\begin{solution}
  \textbf{Output events} stay inside the component tree (one parent listens). \textbf{Pub--sub} suits cross-cutting notifications: cart updated, theme changed, user logged out---sibling or distant components subscribe to the same service channel.
\end{solution}

\question[4]
A simple Angular pub--sub channel in a service might look like:
\begin{choices}
  \choice \lstinline[style=tsinline]|@Component({ bootstrap: [MessageBus] })|
  \CorrectChoice \lstinline[style=tsinline]|private bus = new Subject<string>(); publish(msg) { this.bus.next(msg); } onMessage() { return this.bus.asObservable(); }|
  \choice \lstinline[style=tsinline]|new HttpClient()| inside every template
  \choice \lstinline[style=htmlexam]!{{ bus | date }}! only
\end{choices}
\begin{solution}
  Expose the stream as an \lstinline[style=tsinline]|Observable| (\lstinline[style=tsinline]|asObservable()|) so callers cannot call \lstinline[style=tsinline]|next()| from outside the service. Register the service with \lstinline[style=tsinline]|providedIn: 'root'| for an app-wide bus.
\end{solution}

\question[4]
Main benefits of pub--sub in large Angular apps include:
\begin{choices}
  \choice Eliminating the need for services and DI entirely
  \CorrectChoice Looser coupling, easier addition of new listeners, and centralized event channels
  \choice Automatic SQL generation from templates
  \choice Replacing all route guards
\end{choices}
\begin{solution}
  New features can subscribe without editing the publisher. Trade-off: indirection---too many global buses make data flow hard to trace; prefer explicit services, state stores, or router state when a single source of truth is clearer.
\end{solution}

\question[4]
Which practice helps avoid memory leaks in pub--sub?
\begin{choices}
  \choice Never unsubscribe from any Observable
  \CorrectChoice Unsubscribe in \lstinline[style=tsinline]|ngOnDestroy| (or use \lstinline[style=tsinline]|takeUntil(destroy$)| / \lstinline[style=htmlexam]|async| pipe) when the component is destroyed
  \choice Call \lstinline[style=tsinline]|bootstrapModule| on every message
  \choice Store subscriptions in \texttt{styleUrls}
\end{choices}
\begin{solution}
  Long-lived root services may keep a hot \lstinline[style=tsinline]|Subject| alive; \textbf{component} subscriptions must end when the view goes away. Patterns: \lstinline[style=tsinline]|takeUntil|, manual \lstinline[style=tsinline]|subscription.unsubscribe()|, or template \lstinline[style=htmlexam]|async| pipe.
\end{solution}

\question[4]
Compared to passing callbacks through many component layers, pub--sub via a service:
\begin{choices}
  \choice Always requires \lstinline[style=htmlexam]|*ngSwitch| on the root
  \CorrectChoice Avoids ``prop drilling'' by letting distant components hear the same event
  \choice Prevents use of \lstinline[style=tsinline]|@Input|
  \choice Replaces \texttt{HttpClient} with \texttt{XMLHttpRequest} only
\end{choices}
\begin{solution}
  \textbf{Input/output} chains are explicit and good for direct parent--child data. \textbf{Pub--sub} helps when many unrelated parts of the UI need the same signal---at the cost of making the flow less visible in templates.
\end{solution}

\newpage

\uplevel{\par\textbf{Part 25: RxJS observables and subjects}\par}

\question[4]
An RxJS \textbf{Observable} is best described as:
\begin{choices}
  \choice A single static JSON file in \texttt{assets/}
  \CorrectChoice A lazy stream that can emit zero or more values over time and may complete or error
  \choice A CSS encapsulation mode
  \choice The same as a route \texttt{path} string
\end{choices}
\begin{solution}
  Observers \texttt{subscribe} with handlers for \textbf{next}, \textbf{error}, and \textbf{complete}. \texttt{HttpClient}, DOM events, intervals, and user input are commonly modeled as Observables in Angular.
\end{solution}

\question[4]
Calling \lstinline[style=tsinline]|observable.subscribe({ next: v => ..., error: e => ... })|:
\begin{choices}
  \choice Sends HTTP headers without creating a connection
  \CorrectChoice Registers handlers; for cold streams (such as HTTP), work often begins at subscribe time
  \choice Immediately deletes the Observable from memory always
  \choice Replaces \lstinline[style=tsinline]|@NgModule| imports
\end{choices}
\begin{solution}
  Without a subscriber, many Observables do nothing. \textbf{Cold} Observables (typical HTTP) run per subscription; \textbf{hot} sources (Subjects, DOM events) may emit to current subscribers only.
\end{solution}

\question[4]
The RxJS operator \lstinline[style=tsinline]|map(x => x * 2)|:
\begin{choices}
  \choice Unsubscribes all components automatically
  \CorrectChoice Transforms each emitted value and passes the result downstream
  \choice Replaces \lstinline[style=tsinline]|CanActivate| guards
  \choice Compiles SCSS to CSS
\end{choices}
\begin{solution}
  Operators return new Observables (pure transformations). Examples: \lstinline[style=tsinline]|filter|, \lstinline[style=tsinline]|tap| (side effects), \lstinline[style=tsinline]|switchMap| (switch to inner Observable---common for search boxes), \lstinline[style=tsinline]|catchError|.
\end{solution}

\question[4]
A \textbf{Subject} differs from a plain Observable factory because a Subject:
\begin{choices}
  \choice Cannot call \lstinline[style=tsinline]|next()|
  \CorrectChoice Is both an Observable and an observer---multicasts values to current subscribers
  \choice Only works in Karma tests
  \choice Must be declared in \texttt{bootstrap}
\end{choices}
\begin{solution}
  Use \lstinline[style=tsinline]|new Subject<T>()| when \emph{you} push values (\lstinline[style=tsinline]|subject.next(value)|). \lstinline[style=tsinline]|BehaviorSubject<T>| also stores the \textbf{latest} value for new subscribers (\lstinline[style=tsinline]|getValue()| / initial seed).
\end{solution}

\question[4]
In a template, \lstinline[style=htmlexam]!{{ heroes$ | async }}! avoids manual subscribe because:
\begin{choices}
  \choice The \lstinline[style=htmlexam]|async| pipe compiles TypeScript to JavaScript at build time
  \CorrectChoice The \lstinline[style=htmlexam]|async| pipe subscribes and unwraps the latest emitted value (and cleans up on destroy)
  \choice \lstinline[style=htmlexam]|async| converts Observables into route paths only
  \choice \lstinline[style=htmlexam]|async| replaces \texttt{HttpClientModule}
\end{choices}
\begin{solution}
  Define \lstinline[style=tsinline]|heroes$ = this.service.getHeroes();| in the class, then bind \lstinline[style=htmlexam]!{{ heroes$ | async }}! in HTML. The pipe handles subscribe/unsubscribe for that binding.
\end{solution}

\question[4]
Which pair is appropriate for app-wide state many components read?
\begin{choices}
  \choice \lstinline[style=tsinline]|EventEmitter| in a leaf component only, with no service
  \CorrectChoice \lstinline[style=tsinline]|BehaviorSubject<User | null>| in an injectable service, exposing \lstinline[style=tsinline]|asObservable()| to templates
  \choice \lstinline[style=htmlexam]|*ngIf| on the router outlet only
  \choice Hard-coded globals on \lstinline[style=tsinline]|window|
\end{choices}
\begin{solution}
  \lstinline[style=tsinline]|BehaviorSubject| gives new subscribers the current user immediately. Components call \lstinline[style=tsinline]|userService.setUser(u)| (wrapping \lstinline[style=tsinline]|next|) instead of mutating a shared object from many places without a pattern.
\end{solution}

\uplevel{\par\textbf{Part 26: Karma and Jasmine testing}\par}

\question[4]
In the default Angular CLI setup, unit tests are run with:
\begin{choices}
  \choice \texttt{ng serve} only
  \CorrectChoice \texttt{ng test} (Karma test runner + Jasmine assertions)
  \choice \texttt{ng build --prod} only
  \choice \texttt{npm start} without a test configuration
\end{choices}
\begin{solution}
  \texttt{ng test} launches \textbf{Karma}, which loads the app in a browser (or headless Chrome), executes \texttt{*.spec.ts} files, and reports results. \texttt{ng e2e} (when configured) covers end-to-end tests separately.
\end{solution}

\question[4]
In Jasmine, \lstinline[style=tsinline]|describe('HeroService', () => { ... })| defines:
\begin{choices}
  \choice A production route path
  \CorrectChoice A test suite grouping related \lstinline[style=tsinline]|it| examples
  \choice An Angular \lstinline[style=tsinline]|@NgModule|
  \choice A CSS encapsulation boundary
\end{choices}
\begin{solution}
  Structure: \lstinline[style=tsinline]|describe| (suite) $\rightarrow$ \lstinline[style=tsinline]|beforeEach| (setup) $\rightarrow$ \lstinline[style=tsinline]|it('should ...', () => { expect(...).toBe(...) })| (single spec).
\end{solution}

\question[4]
The matcher \lstinline[style=tsinline]|expect(actual).toEqual(expected)| checks:
\begin{choices}
  \choice Reference identity only (\lstinline[style=tsinline]|===| for objects always)
  \CorrectChoice Deep equality for objects and arrays (value comparison)
  \choice That a route guard returns \texttt{true}
  \choice HTTP status codes only
\end{choices}
\begin{solution}
  Common matchers: \lstinline[style=tsinline]|toBe| (identity), \lstinline[style=tsinline]|toEqual| (deep), \lstinline[style=tsinline]|toBeTruthy|/\lstinline[style=tsinline]|toBeFalsy|, \lstinline[style=tsinline]|toContain|, \lstinline[style=tsinline]|toThrow|. Jasmine also supports spies: \lstinline[style=tsinline]|spyOn(service, 'getHeroes')|.
\end{solution}

\question[4]
To test an Angular component in isolation, you typically use:
\begin{choices}
  \choice \lstinline[style=tsinline]|bootstrapModule| in production \texttt{main.ts} only
  \CorrectChoice \lstinline[style=tsinline]|TestBed.configureTestingModule({ declarations: [Component], ... })| and \lstinline[style=tsinline]|TestBed.createComponent(Component)|
  \choice \lstinline[style=htmlexam]|router-outlet| without imports
  \choice \texttt{HttpClient} against the real internet always
\end{choices}
\begin{solution}
\begin{lstlisting}[style=tsexam]
beforeEach(() => {
  TestBed.configureTestingModule({
    declarations: [HeroListComponent],
    imports: [HttpClientTestingModule],
  });
  fixture = TestBed.createComponent(HeroListComponent);
  component = fixture.componentInstance;
});
\end{lstlisting}
  Call \lstinline[style=tsinline]|fixture.detectChanges()| to run change detection after setup.
\end{solution}

\question[4]
\texttt{HttpClientTestingModule} and \texttt{HttpTestingController} are used to:
\begin{choices}
  \choice Style components with \lstinline[style=htmlexam]|:host|
  \CorrectChoice Mock HTTP requests in tests without hitting a real server
  \choice Generate route guards automatically
  \choice Replace Jasmine with Karma
\end{choices}
\begin{solution}
  Tests call \lstinline[style=tsinline]|httpMock.expectOne('/api/heroes')| and \lstinline[style=tsinline]|req.flush(mockHeroes)| to simulate responses. This keeps service tests fast and deterministic.
\end{solution}

\question[4]
A component spec file generated by the CLI is usually named:
\begin{choices}
  \choice \texttt{hero.component.css}
  \CorrectChoice \texttt{hero.component.spec.ts}
  \choice \texttt{hero.module.routes.ts}
  \choice \texttt{karma.json} beside \texttt{index.html}
\end{choices}
\begin{solution}
  Each schematic component includes a \texttt{.spec.ts} next to \texttt{.component.ts}. Karma configuration (often \texttt{karma.conf.js}) tells the runner which files to load; Jasmine provides the assertion API inside those specs.
\end{solution}

\newpage

\uplevel{\par\textbf{Part 27: Routing and navigation}\par}

\question[4]
Angular \textbf{routing} enables:
\begin{choices}
  \choice Replacing TypeScript with HTML only
  \CorrectChoice Mapping URL paths to views (components) in a SPA without full page reloads
  \choice Disabling \texttt{HttpClient} for all apps
  \choice Global CSS without components
\end{choices}
\begin{solution}
  The router updates the browser URL and swaps components in a \lstinline[style=htmlexam]|<router-outlet>|. Users bookmark \texttt{/products/42}; the app loads the correct feature view client-side.
\end{solution}

\question[4]
The root route configuration is typically imported with:
\begin{choices}
  \choice \lstinline[style=tsinline]|RouterModule.forChild(routes)| in \texttt{AppModule} only
  \CorrectChoice \lstinline[style=tsinline]|RouterModule.forRoot(routes)| in the root module (or \lstinline[style=tsinline]|provideRouter(routes)| in standalone apps)
  \choice \lstinline[style=tsinline]|HttpClientModule.forRoot| instead of routes
  \choice \texttt{declarations: [Router]}
\end{choices}
\begin{solution}
  \textbf{forRoot} registers app-wide router providers once. Feature modules use \textbf{forChild} to add routes without duplicating the singleton router setup.
\end{solution}

\question[4]
Which template link navigates to the \texttt{/heroes} route?
\begin{choices}
  \choice \lstinline[style=htmlexam]|<a href="/heroes">Heroes</a>| only (always full page reload in every setup)
  \CorrectChoice \lstinline[style=htmlexam]|<a routerLink="/heroes">Heroes</a>| (with \lstinline[style=htmlexam]|RouterModule| imported)
  \choice \lstinline[style=htmlexam]|<a ngModel="/heroes">Heroes</a>|
  \choice \lstinline[style=htmlexam]|<a *ngFor="/heroes">Heroes</a>|
\end{choices}
\begin{solution}
  \lstinline[style=htmlexam]|routerLink| (and \lstinline[style=htmlexam]|routerLinkActive| for active CSS) delegates navigation to the Angular router. Plain \lstinline[style=htmlexam]|href| may reload the entire document unless configured for hash routing or server fallback.
\end{solution}

\question[4]
The \lstinline[style=htmlexam]|<router-outlet>| element:
\begin{choices}
  \choice Holds the entire \texttt{package.json} file
  \CorrectChoice Is the placeholder where the router inserts the active route's component
  \choice Replaces \lstinline[style=tsinline]|HttpClient|
  \choice Must appear only in \texttt{main.ts}
\end{choices}
\begin{solution}
  Parent layouts keep chrome (nav bar); child routes render inside the outlet. Nested routes use multiple outlets or child \lstinline[style=htmlexam]|<router-outlet>| elements in feature components.
\end{solution}

\question[4]
Programmatic navigation from TypeScript uses:
\begin{choices}
  \choice \lstinline[style=tsinline]|EventEmitter.navigate()|
  \CorrectChoice \lstinline[style=tsinline]|this.router.navigate(['/heroes', id])| or \lstinline[style=tsinline]|navigateByUrl('/heroes')| after injecting \lstinline[style=tsinline]|Router|
  \choice \lstinline[style=htmlexam]|{{ routerLink }}| interpolation only
  \choice \lstinline[style=tsinline]|bootstrapModule(['/heroes'])|
\end{choices}
\begin{solution}
  Inject \lstinline[style=tsinline]|Router| from \lstinline[style=tsinline]|@angular/router|. Use \lstinline[style=tsinline]|ActivatedRoute| to read params (\lstinline[style=tsinline]|route.snapshot.paramMap.get('id')| or \lstinline[style=tsinline]|paramMap$| Observable).
\end{solution}

\question[4]
A route definition \lstinline[style=tsinline]|{ path: '**', redirectTo: '/home' }| typically:
\begin{choices}
  \choice Bootstraps \texttt{AppComponent}
  \CorrectChoice Acts as a wildcard fallback when no other path matches
  \choice Registers \texttt{HttpClient} providers
  \choice Runs Karma tests automatically
\end{choices}
\begin{solution}
  Order matters: specific paths first, wildcard last. Other common entries: \lstinline[style=tsinline]|{ path: '', redirectTo: 'dashboard', pathMatch: 'full' }|, \lstinline[style=tsinline]|loadChildren| for lazy-loaded feature modules, and \lstinline[style=tsinline]|canActivate| guards (see Part~21).
\end{solution}

\question[4]
Lazy loading a feature module means:
\begin{choices}
  \choice Downloading the entire app in \texttt{index.html} up front always
  \CorrectChoice Loading a feature's module and routes only when the user navigates to that path
  \choice Disabling the router outlet
  \choice Using \lstinline[style=tsinline]|EventEmitter| instead of routes
\end{choices}
\begin{solution}
  Example: \lstinline[style=tsinline]|{ path: 'admin', loadChildren: () => import('./admin/admin.module').then(m => m.AdminModule) }|. Smaller initial bundle; faster first paint. Requires feature routing with \lstinline[style=tsinline]|RouterModule.forChild|.
\end{solution}

\endgradingrange{csawebangular}
